\documentclass[coursework]{SCWorks}

\renewcommand{\baselinestretch}{1.5}

\usepackage{preamble}

\begin{document}

\chair{информатики и программирования}

\title{Разработка интерактивной аналитической панели для визуализации и анализа данных рынка недвижимости на основе данных платформы Яндекс.Недвижимость}

\course{3}
\group{341}

\napravlenie{02.03.03 "--- Математическое обеспечение и администрирование информационных систем}

\author{Перфилова Арсения Сергеевича}

\chtitle{доцент, к.\,ф.-м.\,н.}
\chname{М.\,В.\,Огнева}

\satitle{ст.~преп. каф.~ИиП}
\saname{А.\,А.\,Казачкова}

\date{2026}

\begin{spacing}{1.2}
    \maketitle
\end{spacing}

\secNumbering

\tableofcontents

\intro

На агрегаторах вроде Яндекс.Недвижимость каждый день публикуются тысячи объявлений о продаже квартир: цена, площадь, этаж, тип дома, адрес. По таким данным изучают рынок недвижимости и ищут закономерности в ценообразовании~\cite{soboleva2023,belyaev2015}. При этом сам сайт для анализа неудобен: фильтры не покрывают все нужные комбинации условий, нельзя быстро сравнить два сегмента рынка и посмотреть выборку на карте в одном месте.

Данные ещё и непросто достать. Публичного API у площадки нет, карточки подгружаются через JavaScript, работает защита от автоматического сбора. Обычно приходится отдельно писать парсер, отдельно чистить таблицу и только потом строить графики~\cite{nesterov2024,selenium}. Статические отчёты в Excel или ноутбуке решают часть задач, но не дают интерактивности: при смене фильтра всё нужно пересчитывать заново~\cite{filatov2013,ahtimirov2024}.

В этой работе собран полный цикл в одном приложении: Python-пайплайн забирает объявления с Яндекс.Недвижимость, сервер хранит готовые датасеты по городам, а в браузере открывается дашборд с фильтрами, сводными показателями, графиками и картой. Подход близок к тому, как устроены другие веб-системы для анализа данных~\cite{filimonova2024,sultanov2022}.

Объект работы~--- объявления о продаже квартир на Яндекс.Недвижимость. Предмет~--- способы автоматизированного сбора этих данных, их подготовки и показа в интерактивном веб-интерфейсе.

Цель~--- разработать такую систему для анализа рынка недвижимости: парсер на Python, серверную часть с хранением данных по городам и одностраничное веб-приложение с разными правами доступа.

Чтобы этого добиться, были поставлены задачи:
\begin{enumerate}
    \item Разобраться, как устроен сайт Яндекс.Недвижимость, и выбрать схему сбора данных через Selenium с этапами «ссылки~$\rightarrow$ карточки~$\rightarrow$ подготовка CSV».
    \item Написать парсер с настройками под разные города, отслеживанием прогресса и возможностью продолжить работу после остановки или капчи.
    \item Спроектировать приложение: React и TypeScript на клиенте, Express и SQLite на сервере, общие типы и фильтры в каталоге \texttt{shared/}.
    \item Сделать дашборд~--- фильтры, блок с медианами и числом объектов, настраиваемые графики, карту, таблицу объявлений и сравнение двух наборов фильтров.
    \item Поддержать загрузку данных с сервера (Москва, Саратов) и из своего CSV, экспорт результатов, подстановку адреса по координатам с кэшем.
    \item Разделить пользователей по ролям: просмотр, полный анализ, администрирование датасета и парсера, отдельная панель для руководителя.
    \item Проверить систему на реальных данных и убедиться, что ею можно пользоваться для разбора локального рынка.
\end{enumerate}

\section{Анализ предметной области и теоретические основы}

\subsection{Рынок недвижимости как объект компьютерного анализа}

Рынок жилой недвижимости относится к сферам, где одновременно действуют экономические закономерности, пространственная неоднородность и большой объём разрозненной информации. Для покупателя важны цена, район и состояние дома; для исследователя --- структура предложения, динамика цен, связь характеристик объекта с его стоимостью. Традиционные источники --- отчёты Росреестра, выборочные обследования, публикации аналитических агентств --- дают агрегированную картину, но редко позволяют оперативно построить произвольный срез по микропризнакам: этажности, типу дома, году постройки, удалённости от центра.

С распространением интернет-площадок ситуация изменилась. Агрегаторы объявлений ежедневно публикуют тысячи предложений о продаже квартир, и каждое объявление содержит набор атрибутов: стоимость, площадь, число комнат, адрес или координаты, параметры здания. Как показывают исследования, компьютерный анализ таких массивов помогает выявлять факторы ценообразования и особенности локальных сегментов рынка~\cite{soboleva2023,belyaev2015}. При этом сами площадки ориентированы на поиск жилья, а не на статистику: встроенные фильтры ограничены, сравнение нескольких сегментов рынка неудобно, географическая картина представлена фрагментарно.

Отсюда следует постановка задачи, характерная для прикладной информатики: необходимо не только получить данные из веб-источника, но и привести их к виду, пригодному для анализа, а затем обеспечить наглядное и гибкое исследование --- через специализированные программные средства визуализации и фильтрации~\cite{filatov2013,ahtimirov2024}. Далее в разделе рассматриваются теоретические аспекты каждого из этих этапов: от способов сбора информации в сети до принципов построения интерактивных аналитических панелей.

\subsection{Классификация источников данных о недвижимости}

С точки зрения анализа данных источники можно разделить на первичные и вторичные. К первичным относятся сведения, полученные непосредственно с площадок размещения объявлений или через обращение к ним программными средствами. К вторичным --- уже обработанные датасеты, отчёты и индексы, подготовленные третьими сторонами.

\textbf{Официальная статистика и реестры} дают высокую достоверность по сделкам и правам, но обновляются с задержкой и не всегда содержат детальные характеристики объекта в момент выставления на продажу.

\textbf{Коммерческие агрегаторы} (классифайды, порталы недвижимости) отражают текущее предложение: актуальные цены «запроса», а не факт сделки. Данные богаты атрибутами, но могут содержать дубликаты, устаревшие объявления и ошибки ввода со стороны продавца или агента.

\textbf{Открытые программные интерфейсы (API)} встречаются у крупных сервисов, однако для многих агрегаторов публичный API либо отсутствует, либо ограничен партнёрами. В таких случаях исследователь вынужден обращаться к методам извлечения данных из пользовательского веб-интерфейса.

\textbf{Пользовательские выгрузки} --- CSV, Excel, выгрузки из CRM агентств --- удобны для локального анализа, но плохо масштабируются и не гарантируют единообразие схемы полей.

Для задач мониторинга рынка в режиме, близком к реальному времени, наибольший интерес представляют агрегаторы объявлений. Их данные носят полуструктурированный характер: часть полей формализована (цена, площадь), часть --- текстовое описание (ремонт, инфраструктура), часть --- географические координаты или адрес в свободной форме. Это накладывает требования на этапы очистки, нормализации и контроля качества, рассмотренные в подразделе~1.4.

\subsection{Обзор методов сбора и анализа данных в веб-пространстве}

Сбор данных из интернета --- самостоятельное направление, на стыке программирования, сетевых технологий и прикладной статистики. В обзорной литературе и учебных курсах по анализу данных экосистема Python рассматривается как стандартный инструмент для последующей обработки табличных массивов~\cite{mckinney2023,vanderplas2023,coursera-python}. Библиотека pandas обеспечивает представление данных в виде таблицы (DataFrame), операции фильтрации, группировки и агрегирования; NumPy --- эффективные числовые вычисления над столбцами. Однако до этапа загрузки в DataFrame необходимо решить задачу \emph{извлечения} информации из HTML-страниц или API.

Веб-страница с точки зрения программы --- документ разметки HTML, дополненный каскадными таблицами стилей и скриптами JavaScript. Совокупность элементов страницы образует объектную модель документа (DOM). Аналитически полезные сведения могут находиться в текстовых узлах, атрибутах тегов, встроенных JSON-блоках (в том числе разметке schema.org) или подгружаться асинхронно после действий пользователя.

Различают два принципиально разных подхода к получению содержимого.

\textbf{Статический доступ} предполагает, что полный HTML уже содержит искомые данные. Достаточно выполнить HTTP-запрос, сохранить ответ и разобрать разметку парсером. Классическим инструментом для навигации по дереву тегов служит библиотека Beautiful Soup~\cite{beautifulsoup}: она позволяет находить элементы по имени, классу, идентификатору и извлекать текст или атрибуты.

\textbf{Динамический доступ} необходим, когда содержимое формируется на стороне клиента. Типичны одностраничные приложения (SPA): первичный HTML --- «оболочка», а объявления подгружаются запросами к внутренним API и вставляются в DOM скриптами. В такой ситуации простой HTTP-клиент получает пустой или неполный документ. Требуются средства автоматизации браузера --- в первую очередь Selenium WebDriver~\cite{selenium}, управляющий реальным или headless-браузером: открытие URL, прокрутка, клики, ожидание появления элементов, чтение DOM после выполнения JavaScript.

Сравнение подходов сведено в таблице~\ref{tab:parse-methods}.

\begin{table}[htbp]
\centering
\begin{tabular}{|p{3.2cm}|p{4.8cm}|p{4.8cm}|}
\hline
\textbf{Критерий} & \textbf{Статический парсинг} & \textbf{Автоматизация браузера} \\
\hline
Тип сайта & Серверный рендеринг HTML & SPA, AJAX, бесконечная прокрутка \\
\hline
Скорость & Выше & Ниже (накладные расходы браузера) \\
\hline
Устойчивость к JS & Нет & Да \\
\hline
Ресурсоёмкость & Низкая & Высокая \\
\hline
Типичные инструменты & HTTP-клиент, Beautiful Soup & Selenium, Playwright \\
\hline
\end{tabular}
\caption{Сравнение подходов к извлечению данных из веб-страниц}
\label{tab:parse-methods}
\end{table}

Для агрегаторов недвижимости всё чаще актуален второй подход или комбинация: браузер обеспечивает навигацию и подгрузку, после чего данные считываются из DOM или из перехваченных JSON-ответов.

После сбора начинается этап анализа. В учебниках по Data Science выделяют разведочный анализ данных (Exploratory Data Analysis, EDA) --- систематическое изучение распределений, пропусков, выбросов и связей между переменными до построения прогнозных моделей~\cite{demidova2022,grus2021}. Для описательной аналитики рынка недвижимости EDA нередко является конечной целью: исследователю нужны гистограммы цен, карта предложений, сравнение сегментов, а не обязательно модель машинного обучения. Вопросы интерпретируемости и участия эксперта в аналитическом процессе подробно обсуждаются в литературе по машинному обучению с участием человека~\cite{monarch2022}.

Для регулярного мониторинга рынка и совместного доступа к одним и тем же данным разовых скриптов недостаточно: целесообразно веб-приложение с интерактивной панелью~\cite{filimonova2024,sultanov2022}.

\subsection{Технологии парсинга и ограничения веб-источников}

\subsubsection{Структура HTML и локализация данных}

Проектирование парсера начинается с анализа структуры целевых страниц. HTML представляет иерархию вложенных элементов; CSS-классы и XPath-выражения служат для однозначного указания нужного фрагмента. На странице списка объявлений обычно повторяется один и тот же шаблон карточки; на странице объ объявления --- блоки «основные параметры», «о доме», «расположение». Разработчик фиксирует соответствие между полями будущей таблицы и фрагментами DOM.

Помимо видимого HTML данные могут дублироваться во встроенных JSON-LD-блоках (формат linked data для поисковых систем). Использование нескольких источников внутри одной страницы повышает полноту извлечения: если текстовое поле пусто, значение может быть получено из структурированной разметки.

\subsubsection{Поэтапный pipeline сбора}

На практике сбор с крупных каталогов редко сводится к одному проходу по сайту. Логично выделяют стадии:
\begin{enumerate}
    \item \textbf{сбор ссылок} --- обход страниц поиска с разными параметрами (район, число комнат), накопление URL карточек без детального разбора;
    \item \textbf{сбор деталей} --- последовательный обход списка ссылок, извлечение полного набора атрибутов каждого объявления;
    \item \textbf{нормализация} --- приведение «сырой» таблицы к единой схеме, расчёт производных показателей, удаление дубликатов;
    \item \textbf{контроль качества} --- проверка полноты ключевых полей, доли пропусков, географической правдоподобности координат.
\end{enumerate}
Такое разделение упрощает возобновление после сбоя, повторный прогон только одной стадии и параллельную настройку параметров обхода для разных регионов.

\subsubsection{Защита от автоматизированного доступа}

Владельцы коммерческих площадок ограничивают интенсивность автоматических обращений. К распространённым мерам относятся:
\begin{itemize}
    \item \textbf{ограничение частоты запросов} --- временные блокировки IP при превышении порога;
    \item \textbf{CAPTCHA} --- задачи, требующие участия человека;
    \item \textbf{проверка заголовков и отпечатка браузера} --- отличие скрипта от обычного пользователя;
    \item \textbf{динамическая структура разметки} --- изменение классов и путей DOM, затрудняющее поддержку парсера.
\end{itemize}
С теоретической точки зрения парсинг превращается в задачу устойчивого извлечения при неполной информации о внутреннем устройстве сайта и при активном противодействии~\cite{nesterov2024}. Рекомендуется варьировать паузы между запросами, имитировать естественную прокрутку, корректно обрабатывать сценарии с CAPTCHA (пауза, ручное решение, продолжение). Эти меры не гарантируют бессрочной работоспособности парсера, но снижают риск блокировки.

\subsubsection{Альтернатива парсингу --- программный интерфейс}

Если площадка предоставляет документированный API с авторизацией, предпочтительно использовать его: данные приходят в JSON или XML, схема полей известна, соблюдаются лимиты использования. Отсутствие открытого API вынуждает прибегать к парсингу пользовательского интерфейса, что повышает стоимость сопровождения при каждом изменении вёрстки сайта.

\subsection{Методология разведочного анализа и предобработки данных}

Сырые данные с веб-источника редко пригодны для немедленной визуализации. Ошибки ввода, пропуски, несогласованные форматы чисел и дат, текстовые поля вместо категорий --- типичные проблемы~\cite{mckinney2023}. Предобработка и EDA направлены на превращение такого массива в «чистую» таблицу с понятной семантикой столбцов.

\subsubsection{Пропущенные значения}

Механизмы появления пропусков классифицируют как полностью случайные (MCAR), случайные (MAR) и неслучайные (MNAR)~\cite{demidova2022}. Выбор стратегии зависит от доли пропусков и от того, связана ли их природа с целевым показателем (например, с ценой).

Возможные действия:
\begin{itemize}
    \item \textbf{удаление строк} --- оправдано при малой доле потерь в некритичных полях или когда пропуск означает отсутствие объекта на рынке;
    \item \textbf{удаление столбцов} --- если признак заполнен менее чем на заданный процент;
    \item \textbf{импутация} --- заполнение средним, медианой, модой или константой «не указано» для категорий;
    \item \textbf{маркировка} --- отдельный индикатор «значение отсутствовало», если пропуск сам несёт информацию.
\end{itemize}
Для цены и площади квартиры удаление неполных записей часто предпочтительнее заполнения средним, чтобы не искажать распределение.

\subsubsection{Выбросы и аномалии}

Выброс --- наблюдение, существенно отклоняющееся от основной массы данных. На рынке недвижимости выбросами могут быть опечатки (лишний ноль в цене), тестовые объявления или элитные объекты, не представительные для массового сегмента.

Классический метод --- межквартильный размах (IQR). Для переменной $X$ вычисляют первый $Q_1$ и третий $Q_3$ квартили, $IQR = Q_3 - Q_1$. Подозрительными считают значения вне интервала
\begin{equation}
    [Q_1 - 1{,}5 \cdot IQR;\; Q_3 + 1{,}5 \cdot IQR].
    \label{eq:iqr}
\end{equation}
Альтернатива --- z-score: отклонение более чем на $3\sigma$ от среднего. Медиана и квартили устойчивее к выбросам, чем среднее арифметическое~\cite{vanderplas2023,grus2021}.

Решение об удалении или сохранении выброса должно опираться на предметную область: уникальный пентхаус и ошибка парсера требуют разной реакции.

\subsubsection{Распределения и описательная статистика}

Стоимость жилья на вторичном рынке часто имеет правостороннее асимметричное распределение: много объектов в «среднем» сегменте и длинный «хвост» дорогих предложений. Логнормальная модель нередко лучше описывает цены, чем нормальная. Для описания центра распределения в таких случаях предпочитают \textbf{медиану}; среднее арифметическое чувствительно к дорогим объектам. \textbf{Мода} полезна для дискретных признаков (число комнат).

Понимание формы распределения влияет на выбор визуализации (шаг гистограммы, шкала осей) и на интерпретацию трендов.

\subsubsection{Корреляция и связи между признаками}

Между площадью и ценой часто наблюдается положительная связь. Коэффициент корреляции Пирсона $r$ измеряет линейную зависимость двух количественных переменных:
\begin{equation}
    r = \frac{\sum (x_i - \bar{x})(y_i - \bar{y})}{\sqrt{\sum (x_i - \bar{x})^2 \sum (y_i - \bar{y})^2}}.
    \label{eq:pearson}
\end{equation}
Значения $r$ близкие к $+1$ или $-1$ указывают на сильную линейную связь; близкие к нулю --- на её отсутствие. Корреляция не доказывает причинно-следственную связь и плохо улавливает нелинейные зависимости, поэтому её дополняют диаграммами рассеяния~\cite{monarch2022}.

\subsubsection{Инженерия признаков}

Исходные поля объявления можно обогатить производными признаками, облегчающими анализ:
\begin{itemize}
    \item \textbf{удельная цена} --- отношение стоимости к площади (руб./м²);
    \item \textbf{возраст здания} --- разность между текущим годом и годом постройки;
    \item \textbf{относительный этаж} --- отношение этажа к этажности дома;
    \item \textbf{расстояние до условного центра города} --- по формуле haversine из географических координат;
    \item \textbf{бинарные индикаторы} --- наличие балкона, лифта, парковки и т.п.
\end{itemize}

Категориальные переменные (тип дома: панельный, кирпичный, монолитный) для алгоритмов фильтрации и некоторых моделей преобразуют в числовой вид. \textbf{One-hot encoding} создаёт отдельный столбец для каждой категории со значениями 0 или 1. Альтернатива --- порядковое кодирование, если категории имеют естественный порядок (например, состояние ремонта).

\subsubsection{Дедупликация и слияние записей}

Одно и то же объявление может быть собрано повторно при разных запусках парсера. Дубликаты искажают частоты и средние показатели. Идентификация ведётся по стабильному ключу --- URL объявления, внешнему идентификатору на площадке или комбинации адреса, площади и цены. При слиянии двух версий одной записи из разных выгрузок логично объединять непустые поля: если в первой версии не было координат, а во второй они появились, результирующая строка должна их сохранить.

\subsubsection{Контроль качества данных}

Качество данных оценивают по нескольким измерениям~\cite{demidova2022}:
\begin{itemize}
    \item \textbf{полнота} --- доля заполненных значений в ключевых столбцах;
    \item \textbf{уникальность} --- отсутствие нежелательных дубликатов;
    \item \textbf{валидность} --- соответствие форматов и допустимых диапазонов (цена ${>}0$, площадь в разумных пределах);
    \item \textbf{согласованность} --- непротиворечивость связанных полей (этаж не больше этажности дома);
    \item \textbf{актуальность} --- давность сбора относительно момента анализа.
\end{itemize}
Формализованные пороги (например, не менее 75\% заполненности координат, не более 15\% пропусков в цене) позволяют автоматически отклонять низкокачественные выгрузки до этапа визуализации.

\subsection{Теоретические основы визуализации данных}

Визуализация переводит числовую таблицу в форму, удобную для зрительного восприятия. Человек быстрее замечает тренды, кластеры и аномалии на графике, чем при просмотре тысяч строк~\cite{vanderplas2023}. Цель визуализации в аналитике недвижимости --- не украшение отчёта, а снижение когнитивной нагрузки при исследовании многомерных данных.

\subsubsection{Принципы эффективного графического представления}

При выборе типа диаграммы учитывают вид переменных:
\begin{itemize}
    \item одна количественная --- гистограмма, box plot, плотность;
    \item одна категориальная --- столбчатая диаграмма частот;
    \item две количественные --- диаграмма рассеяния;
    \item количественная и категориальная --- grouped bar chart, box plot по группам;
    \item географические координаты --- точечная карта, тепловая карта плотности.
\end{itemize}
Нарушение соответствия типа графика типу данных (например, гистограмма для номинальных категорий без порядка) ведёт к неверной интерпретации.

\subsubsection{Гистограммы и биннинг}

Гистограмма делит диапазон значений на интервалы (bins) и показывает число наблюдений в каждом. Ширина интервала влияет на вид графика: слишком узкие bins дают шум, слишком широкие скрывают мультимодальность.

Правило Стерджеса предлагает число интервалов $k = \lceil \log_2 n \rceil + 1$ для выборки из $n$ точек. Правило Фридмана--Диакониса учитывает размах данных и межквартильный размах:
\begin{equation}
    h = 2 \cdot IQR \cdot n^{-1/3},
    \label{eq:fd}
\end{equation}
где $h$ --- ширина bin~\cite{grus2021}. На практике допустима интерактивная подстройка шага пользователем.

\subsubsection{Диаграммы рассеяния и проблема наслоения}

При построении scatter-графика «цена--площадь» для большой выборки точки накладываются друг на друга (overplotting), и плотность участков становится неразличимой. Приёмы борьбы: полупрозрачность маркеров (alpha blending), субдискретизация (random sampling), контурные линии плотности, шестигранная бининг-диаграмма (hexbin)~\cite{mckinney2023}.

Добавление линии тренда (например, методом наименьших квадратов) облегчает оценку общего направления связи, но не должно подменять анализ остатков и выбросов.

\subsubsection{Кодирование категорий цветом и формой}

Для различения групп (тип дома, район) используются визуальные каналы: цвет, форма маркера, тип линии. Рекомендуется не перегружать один график большим числом категорий --- зрительное восприятие уверенно различает ограниченное количество оттенков. Категориальные шкалы предпочтительнее раскрашивать качественными палитрами без ложного порядка «больше-меньше»~\cite{monarch2022}.

\subsection{Интерактивные аналитические панели (дашборды)}

\subsubsection{Отличие дашборда от статического отчёта}

Статический отчёт (PDF, слайды, фиксированная таблица Excel) фиксирует один срез данных на момент подготовки. \textbf{Интерактивный дашборд} позволяет пользователю менять фильтры, масштаб и состав виджетов, получая обновлённую картину без повторного построения документа~\cite{filatov2013,filimonova2024}. Связь «действие пользователя $\rightarrow$ пересчёт выборки $\rightarrow$ обновление графиков» лежит в основе современных систем бизнес-аналитики и специализированных исследовательских интерфейсов.

Преимущества интерактивности для анализа недвижимости:
\begin{itemize}
    \item быстрая проверка гипотез («как меняется медиана цены при ограничении радиуса от центра»);
    \item сопоставление сегментов («однокомнатные vs двухкомнатные») на одном экране;
    \item географическая привязка --- карта как дополнение к табличным и графическим представлениям;
    \item доступность для пользователей без навыков программирования.
\end{itemize}

\subsubsection{Архитектура представления данных в дашборде}

Концептуально дашборд опирается на три слоя:
\begin{enumerate}
    \item \textbf{хранилище} --- таблица наблюдений (строки --- объявления, столбцы --- признаки);
    \item \textbf{логика фильтрации} --- правила отбора строк по диапазонам и категориям;
    \item \textbf{слой представления} --- виджеты (KPI, графики, карта, таблица), каждый из которых получает уже отфильтрованную или агрегированную выборку.
\end{enumerate}
Изменение фильтра пересчитывает все зависимые виджеты. При большом объёме данных критична производительность: отбор строк и агрегирование должны выполняться за время, не нарушающее плавность интерфейса (ориентир --- доли секунды на типичной выборке).

\subsubsection{Типовые элементы дашборда для рынка недвижимости}

\textbf{KPI-блок (ключевые показатели)} --- краткие числа по текущей выборке: число объектов, медиана цены, медиана цены за м², медиана площади. Медианы предпочтительнее средних при асимметричных распределениях цен.

\textbf{Панель фильтров} --- элементы ввода для цены, площади, комнатности, года постройки, типа дома, расстояния от центра. Дополнительные условия по произвольным столбцам расширяют гибкость без изменения кода приложения.

\textbf{Набор графиков} --- распределения, сравнения категорий, scatter с трендом. Возможность включать и отключать отдельные графики снижает перегруженность экрана.

\textbf{Карта} --- точечное отображение объектов с кластеризацией при большой плотности; всплывающие подсказки с ценой и адресом.

\textbf{Табличное представление} --- детальный просмотр отдельных объявлений с сортировкой по столбцам и ссылкой на первоисточник.

\textbf{Сравнение сегментов} --- два независимых набора фильтров и параллельное отображение KPI или графиков для каждого.

\subsubsection{Географическая визуализация}

Недвижимость по своей природе связана с географией. Карта дополняет табличный анализ: видны «горячие» зоны предложения, радиальные градиенты цены, изолированные кластеры. Точечные карты подходят для отдельных объявлений; хороплетные карты (раскраска районов по агрегированному показателю) --- для сравнения территорий, если выборка достаточна для устойчивых средних.

При отсутствии текстового адреса, но наличии координат применяют \textbf{обратное геокодирование} (reverse geocoding) --- преобразование пары (широта, долгота) в адрес или название места с помощью внешних геосервисов. Результаты кэшируют, чтобы не обращаться к API повторно для одних и тех же точек.

\subsubsection{Принципы проектирования интерфейса}

При проектировании аналитической панели рекомендуется:
\begin{itemize}
    \item отделить область управления (фильтры) от области результатов (графики, карта);
    \item обеспечить немедленную обратную связь при изменении параметров;
    \item поддерживать адаптивную вёрстку для разных размеров экрана;
    \item предусмотреть светлую и тёмную тему для длительной работы;
    \item не перегружать один экран: при необходимости использовать вкладки или сворачиваемые блоки.
\end{itemize}
Подход, при котором основной экран сразу показывает сводку и виджеты («dashboard first»), распространён в системах мониторинга и исследовательских UI~\cite{filimonova2024}.

\subsection{Обзор программных средств анализа данных}

\subsubsection{Язык Python и библиотеки обработки таблиц}

Python закрепился в качестве универсального языка для анализа данных благодаря библиотекам общего назначения~\cite{habr-bigdata,w3schools-numpy}. Pandas предоставляет структуры DataFrame и Series, операции join, groupby, pivot; NumPy --- векторизованные вычисления. Для задач недвижимости типичны операции: группировка по району, расчёт медианы цены за м², фильтрация по квантилям.

\subsubsection{Среды разработки программных комплексов}

Проект, объединяющий Python-парсер, серверную часть и веб-клиент, удобно вести в интегрированной среде разработки (IDE), поддерживающей несколько языков и запуск скриптов из терминала. Такие среды (Visual Studio Code, JetBrains WebStorm и аналоги) дают подсветку синтаксиса, навигацию по проекту, отладку и единый каталог для всех компонентов системы. Для многопользовательской аналитики и разграничения прав всё равно требуется полноценное веб-приложение, а не локальные скрипты~\cite{ahtimirov2024}.

\subsubsection{Веб-технологии для интерактивной визуализации}

Клиентские библиотеки (D3.js, Chart.js, Recharts и др.) строят графики в браузере, реагируя на изменение данных. Серверная часть может ограничиваться выдачей JSON и авторизацией; фильтрация и отрисовка выполняются на клиенте --- паттерн, характерный для одностраничных приложений~\cite{sultanov2022,nesterov2024}. Выбор конкретного фреймворка определяется требованиями к типизации, экосистеме компонентов и производительности; теоретически любая declarative charting library, интегрируемая с реактивной моделью UI, пригодна для дашборда.

\subsection{Выводы по первому разделу}

В первом разделе рассмотрены теоретические основы, необходимые для построения системы анализа рынка недвижимости на основе веб-данных:
\begin{enumerate}
    \item агрегаторы объявлений --- актуальный, но «шумный» источник; данные требуют очистки и контроля качества;
    \item для динамических сайтов обоснован переход от статического парсинга HTML к автоматизации браузера; сбор целесообразно организовать поэтапно;
    \item предобработка опирается на методы описательной статистики: работа с пропусками, выбросами, производными признаками, кодирование категорий;
    \item интерактивный дашборд даёт гибкость, недостижимую для статического отчёта; его элементы (KPI, фильтры, графики, карта) согласованы с типами анализируемых признаков;
    \item программная реализация может опираться на связку Python (сбор и ETL) и веб-стек (визуализация и доступ пользователей).
\end{enumerate}


\section{Проектирование и архитектура программного комплекса}

В первом разделе рассмотрены теоретические основы сбора, подготовки и визуализации данных о недвижимости. Раздел посвящён проектным решениям: требованиям к системе, составу компонентов, хранению данных, серверному интерфейсу, ролям пользователей и экранам дашборда.

\subsection{Требования к программному комплексу}

Исходя из постановки задачи (см.~введение), сформулированы функциональные и нефункциональные требования.

\textbf{Функциональные требования:}
\begin{enumerate}
    \item автоматизированный сбор объявлений с платформы Яндекс.Недвижимость с возможностью поэтапного выполнения и возобновления после остановки или капчи;
    \item подготовка табличного датасета с единой схемой полей (цена, площадь, комнатность, координаты, параметры дома, категориальные признаки);
    \item хранение серверных снимков рынка по нескольким городам и загрузка пользовательских CSV-файлов;
    \item интерактивный дашборд: фильтры, сводные показатели, настраиваемые графики, карта объектов, сравнение двух сегментов рынка, таблица объявлений;
    \item экспорт отфильтрованных данных и графиков;
    \item обратное геокодирование координат в адрес с кэшированием результатов;
    \item разграничение доступа по ролям: просмотр, расширенная аналитика, администрирование данных и парсера, управление пользователями;
    \item удалённый запуск стадий парсера и контроль прогресса из веб-интерфейса администратора.
\end{enumerate}

\textbf{Нефункциональные требования:}
\begin{itemize}
    \item \emph{отзывчивость интерфейса} --- смена фильтров не должна блокировать ввод; пересчёт показателей выполняется на клиенте без полной перезагрузки страницы;
    \item \emph{согласованность фильтрации} --- одни и те же правила отбора строк применяются в интерфейсе, при серверной пагинации и экспорте;
    \item \emph{масштаб выборки} --- поддержка десятков тысяч объявлений в одном городском снимке (до 50~000 строк);
    \item \emph{развёртывание без отдельного СУБД-сервера} --- для учебного и локального использования достаточно встроенной SQLite и файлового хранения таблиц;
    \item \emph{расширяемость по городам} --- добавление нового рынка сводится к конфигурации парсера и каталогу данных без переписывания клиента.
\end{itemize}

На момент проектирования в систему включены два города --- Москва и Саратов; архитектура рассчитана на добавление других регионов по тому же шаблону.

\subsection{Общая архитектура системы}

Программный комплекс построен по схеме «клиент --- сервер --- внешний источник» с выделенным автономным модулем сбора данных. Логически выделяются четыре части:
\begin{itemize}
    \item \textbf{клиентское веб-приложение} --- одностраничное приложение (SPA) в браузере: авторизация, дашборд, панели администратора;
    \item \textbf{сервер приложений} --- REST API на Node.js: выдача датасетов, авторизация, геокодирование, управление задачами парсера и учётными записями;
    \item \textbf{модуль парсера} --- Python-пайплайн на базе Selenium: обход сайта, извлечение карточек, ETL и контроль качества;
    \item \textbf{общий слой типов и правил} --- каталог с общими для клиента и сервера описаниями городов, фильтров и форматов данных, чтобы логика не расходилась между уровнями.
\end{itemize}

Схема взаимодействия компонентов представлена на рисунке~\ref{fig:architecture}.

\begin{figure}[htbp]
\centering
\begin{tikzpicture}[
    comp/.style={draw, rounded corners=2pt, align=center, minimum height=12mm, inner sep=5pt, font=\small},
    store/.style={comp, fill=gray!10, minimum width=32mm, text width=30mm},
    proc/.style={comp, fill=blue!8, minimum width=34mm, text width=32mm},
    ext/.style={comp, fill=orange!12, dashed, minimum width=34mm, text width=32mm},
    lbl/.style={font=\scriptsize, fill=white, inner sep=1.5pt},
    arr/.style={-{Stealth[length=2.2mm]}, semithick}
]
\node[proc] (client) {Браузер\\React SPA};
\node[proc, below=9mm of client] (api) {Express API\\порт 3001};

\coordinate (rowcenter) at ($(api.south) + (0,-14mm)$);
\node[store] (json) at ([xshift=-42mm]rowcenter) {JSON\\датасеты\\по городам};
\node[store] (sqlite) at (rowcenter) {SQLite\\users, jobs, cache};
\node[proc] (parser) at ([xshift=42mm]rowcenter) {Python-парсер\\Selenium + pandas};

\node[fit=(json)(sqlite)(parser), inner sep=6pt, draw=none] (bottomrow) {};

\node[ext, below=12mm of parser] (yandex) {realty.yandex.ru};
\node[ext, right=24mm of api] (geocoder) {Геосервисы\\Яндекс, OSM};

\draw[arr] (client) -- node[right, lbl] {HTTP /api} (api);
\coordinate (apibranch) at ($(api.south)!0.35!(rowcenter)$);
\draw[arr] (api.south) -- (apibranch);
\draw[arr] (apibranch) -| (json.north);
\draw[arr] (apibranch) -- (sqlite.north);
\draw[arr] (apibranch) -| node[near start, xshift=3mm, lbl, above] {subprocess} (parser.north);
\draw[arr] (parser) -- (yandex);
\draw[arr] (api.east) -- (geocoder.west);
\end{tikzpicture}
\caption{Логическая архитектура программного комплекса}
\label{fig:architecture}
\end{figure}

\textbf{Режим разработки.} Клиент собирается Vite и обслуживается на порту~3000; запросы к \texttt{/api} проксируются на сервер (порт~3001), что даёт единый origin в браузере. В продуктивной сборке статика клиента отдаётся nginx из каталога \texttt{dist/}, API --- тем же Express-приложением за reverse-proxy (см.~раздел~3, подраздел о развёртывании).

\textbf{Разделение ответственности.} Парсер работает как изолированный Python-процесс: сервер только запускает стадии, читает файлы статуса и по завершении импортирует готовый CSV в JSON-хранилище. Визуализация целиком сосредоточена во frontend; сервер не генерирует HTML-страницы аналитики. Личный CSV аналитика хранится в браузере (\texttt{localStorage}) и не попадает на сервер --- это сознательное упрощение для роли «аналитик» без общего файлового хранилища.

\subsection{Проектирование потока данных и пайплайна парсера}

Сбор данных организован поэтапно --- в соответствии с теоретическим обоснованием поэтапного парсинга (раздел~1). Проектный поток приведён на рисунке~\ref{fig:parser-flow}.

\begin{figure}[htbp]
\centering
\begin{tikzpicture}[
    stage/.style={draw, rounded corners, minimum width=2.1cm, minimum height=0.75cm, align=center, font=\small},
    arr/.style={-{Stealth[length=2mm]}}
]
\node[stage] (links) {links};
\node[stage, right=0.55cm of links] (details) {details};
\node[stage, right=0.55cm of details] (etl) {etl};
\node[stage, right=0.55cm of etl] (qa) {qa};
\node[stage, right=0.55cm of qa] (ingest) {ingest};
\draw[arr] (links) -- (details);
\draw[arr] (details) -- (etl);
\draw[arr] (etl) -- (qa);
\draw[arr] (qa) -- (ingest);
\end{tikzpicture}
\caption{Основные стадии пайплайна сбора и подготовки данных}
\label{fig:parser-flow}
\end{figure}

\begin{enumerate}
    \item \textbf{links} --- обход поисковых сегментов города в браузере, накопление URL объявлений;
    \item \textbf{details} --- посещение карточек, извлечение цены, площадей, адреса, координат, параметров здания; поддержка буферизации и возобновления;
    \item \textbf{etl} --- нормализация адресов и координат, расчёт производных признаков (цена за м², расстояние до центра, возраст дома, one-hot кодирование категорий), отсечение заведомо некорректных записей;
    \item \textbf{qa} --- сверка с контрактом схемы: обязательные поля, дубликаты ссылок, выбросы по координатам;
    \item \textbf{ingest} --- загрузка итогового CSV в серверное JSON-хранилище с дедупликацией по идентификатору объявления.
\end{enumerate}

Промежуточные артефакты (списки ссылок, сырые и обработанные таблицы) сохраняются в каталоге данных парсера отдельно для каждого города. Сервер и UI опрашивают JSON-файл runtime-статуса: стадия, процент выполнения, признак капчи --- это позволяет администратору продолжить сбор вручную после прохождения проверки.

Контракт колонок датасета задаёт обязательный минимум (цена, комнаты, площадь, широта, долгота), идентификационные поля (ID, адрес, ссылка) и расширенный набор производных и категориальных признаков (тип дома, ремонт, лифты, парковка и др.). Единый контракт связывает выход ETL, проверку QA и ожидания дашборда.

\subsection{Проектирование модели данных}

Данные системы делятся на \emph{табличные снимки рынка} (основной объём) и \emph{служебные записи} в SQLite (компактные, транзакционные).

\subsubsection{Хранение объявлений}

Каждый город имеет каталог с двумя JSON-файлами:
\begin{itemize}
    \item \texttt{dataset.json} --- массив строк объявлений (\texttt{Record<string, unknown>});
    \item \texttt{dataset-meta.json} --- число строк, дата обновления, история загрузок (batchId, режим append/replace, статистика ingest).
\end{itemize}

Ограничения проектирования: не более 50~000 строк и 20~МБ на одну загрузку. При добавлении партии строки с одинаковым ID объявления сливаются: координаты и ссылки из разных выгрузок дополняют друг друга. Служебные поля (\texttt{\_\_uploadBatchId}, нормализованный тип дома) не экспортируются пользователю.

Ключевые группы столбцов датасета сведены в таблице~\ref{tab:dataset-fields}.

\begin{table}[htbp]
\centering
\small
\setlength{\tabcolsep}{3pt}
\begin{tabular}{|p{0.17\textwidth}|p{0.78\textwidth}|}
\hline
\textbf{Группа} & \textbf{Содержание} \\
\hline
Идентификация & ID объявления, ссылка на карточку, текстовый адрес \\
\hline
Цена и площадь & Цена (руб.), общая и жилая площадь, число комнат, цена за м² \\
\hline
Геопозиция & Широта, долгота, расстояние до центра города (км) \\
\hline
Здание & Год постройки, возраст дома, этаж, этажность, тип дома, лифты \\
\hline
Категории & One-hot признаки: ремонт, санузел, окна, парковка, двор, мебель, техника и др. \\
\hline
\end{tabular}
\caption{Логические группы полей датасета объявлений}
\label{tab:dataset-fields}
\end{table}

\subsubsection{Служебные таблицы SQLite}

В embedded-базе \texttt{app.db} хранятся сущности, для которых нужны транзакции и индексы (таблица~\ref{tab:sqlite}).

\begin{table}[htbp]
\centering
\small
\setlength{\tabcolsep}{3pt}
\begin{tabular}{|L{0.27\textwidth}|L{0.67\textwidth}|}
\hline
\textbf{Таблица} & \textbf{Назначение} \\
\hline
\texttt{users} & Учётные записи: логин, хеш пароля, роль (\texttt{observer}, \texttt{analyst}, \texttt{admin}, \texttt{manager}) \\
\hline
\texttt{parser\_jobs} & Задачи парсера: город, стадия, статус, PID процесса, пути к логу и конфигу, код завершения \\
\hline
\texttt{reverse\_geocode\_cache} & Кэш обратного геокодирования: ключ координат, адрес, срок действия (TTL, по умолчанию 14 суток) \\
\hline
\end{tabular}
\caption{Служебные таблицы SQLite}
\label{tab:sqlite}
\end{table}

Выбор JSON для массива объявлений и SQLite для пользователей и кэша обусловлен балансом простоты развёртывания и достаточной производительности при локальной работе: полный снимок города загружается в память клиента один раз (\texttt{bootstrap}), дальнейшая фильтрация идёт без обращений к серверу на каждый слайдер.

\subsection{Проектирование программного интерфейса сервера}

API построен в REST-стиле; авторизация --- JWT (Bearer-токен, срок жизни 12~часов). Эндпоинты сгруппированы по уровню доступа (таблица~\ref{tab:api}).

\begin{table}[htbp]
\centering
\footnotesize
\setlength{\tabcolsep}{3pt}
\textbf{Эндпоинты для авторизованных пользователей}\\[3pt]
\begin{tabular}{|L{0.34\textwidth}|L{0.11\textwidth}|L{0.49\textwidth}|}
\hline
\textbf{Путь} & \textbf{Метод} & \textbf{Назначение} \\
\hline
\texttt{/api/auth/login} & POST & Аутентификация, выдача JWT \\
\hline
\texttt{/api/me} & GET & Текущий пользователь \\
\hline
\texttt{/api/dataset/bootstrap} & GET & Полная выгрузка строк и сводки по городу \\
\hline
\texttt{/api/summary} & GET & Сводные показатели без всех строк \\
\hline
\texttt{/api/listings} & GET & Пагинация отфильтрованных объявлений \\
\hline
\texttt{/api/geocode/reverse} & GET & Обратное геокодирование одной точки \\
\hline
\end{tabular}

\medskip
\textbf{Эндпоинты analyst и admin}\\[3pt]
\begin{tabular}{|L{0.34\textwidth}|L{0.11\textwidth}|L{0.49\textwidth}|}
\hline
\textbf{Путь} & \textbf{Метод} & \textbf{Назначение} \\
\hline
\texttt{/api/export.csv} & GET & Экспорт CSV (analyst, admin) \\
\hline
\texttt{/api/admin/dataset} & POST/DEL & Загрузка и очистка серверного датасета \\
\hline
\texttt{/api/admin/dataset/}\newline\texttt{ingest-pipeline} & POST & Импорт CSV после ETL \\
\hline
\texttt{/api/admin/parser/jobs} & GET/POST & Список и запуск стадий парсера \\
\hline
\texttt{/api/admin/parser/jobs/}\newline\texttt{:id/resume} & POST & Продолжение после капчи \\
\hline
\end{tabular}

\medskip
\textbf{Эндпоинты manager}\\[3pt]
\begin{tabular}{|L{0.34\textwidth}|L{0.11\textwidth}|L{0.49\textwidth}|}
\hline
\textbf{Путь} & \textbf{Метод} & \textbf{Назначение} \\
\hline
\texttt{/api/manager/markets/overview} & GET & Сводка по всем городам \\
\hline
\texttt{/api/manager/users} & GET/POST & Управление учётными записями \\
\hline
\texttt{/api/manager/users/:id} & PATCH/DEL & Смена роли, пароля, удаление \\
\hline
\end{tabular}
\caption{Основные группы эндпоинтов REST API}
\label{tab:api}
\end{table}

Параметры фильтрации передаются в query string в едином формате, согласованном с клиентским модулем фильтров: диапазоны числовых полей, множественный выбор категорий, дополнительные условия по произвольным столбцам. Это гарантирует, что экспорт и постраничная выдача совпадают с тем, что пользователь видит на дашборде.

\subsubsection{Интеграция API Яндекса}

Для отображения карты объектов и обратного геокодирования координат в адрес система использует два сервиса экосистемы Яндекса: JavaScript API~2.1 на клиенте и API Геокодера на сервере (эндпоинт \texttt{/api/geocode/reverse}). Оба ключа выпускаются в едином кабинете разработчика; на рисунке~\ref{fig:yandex-api} показаны подключённые сервисы и переменные окружения \srce{VITE_YANDEX_MAPS_API_KEY} и \srce{YANDEX_GEOCODER_API_KEY}.

\begin{figure}[htbp]
\centering
\includegraphics[width=0.88\textwidth]{yandex-developer-console}
\caption{Кабинет разработчика Яндекса: ключи JavaScript API и API Геокодера}
\label{fig:yandex-api}
\end{figure}

На скриншоте видно, что для клиентской карты активирован JavaScript API, а для серверного эндпоинта обратного геокодирования --- API Геокодера; разделение ключей по уровням приложения снижает риск несанкционированного использования серверного ключа в браузере. Результаты геокодирования кэшируются в таблице \texttt{reverse\_geocode\_cache} (см.~табл.~\ref{tab:sqlite}).

\subsection{Проектирование ролевой модели}

Система рассчитана на несколько типов пользователей организации, изучающей рынок недвижимости. Матрица возможностей приведена в таблице~\ref{tab:roles}.

\begin{table}[htbp]
\centering
\small
\setlength{\tabcolsep}{3pt}
\begin{tabular}{|p{0.16\textwidth}|p{0.24\textwidth}|p{0.54\textwidth}|}
\hline
\textbf{Роль} & \textbf{Экран} & \textbf{Возможности} \\
\hline
observer & Дашборд & Серверный датасет; KPI, фильтры, фиксированные графики, карта; без таблицы, экспорта и своего CSV \\
\hline
analyst & Дашборд & Всё у observer плюс свой CSV, конструктор графиков, таблица, экспорт CSV/XLSX и PNG/PDF \\
\hline
admin & Дашборд + админ-панели & Управление серверным датасетом, запуск парсера, ingest; без блока «Обзор рынка» \\
\hline
manager & Отдельная консоль & Сводка по рынкам и CRUD пользователей; без доступа к дашборду \\
\hline
\end{tabular}
\caption{Роли пользователей и разграничение доступа}
\label{tab:roles}
\end{table}

Руководитель (\texttt{manager}) не может назначать роль \texttt{manager} другим пользователям --- это защита от эскалации прав. Наблюдатель получает упрощённый список объявлений вместо полной таблицы, чтобы снизить риск массового копирования данных (см.~рис.~\ref{fig:observer-dashboard} в подразделе~2.7.2). Администратор при пустом городе попадает на экран первичной настройки: предложение загрузить CSV или запустить сбор; панели ingest и парсера --- на рис.~\ref{fig:admin-server} и~\ref{fig:admin-parser}. Консоль руководителя --- на рис.~\ref{fig:manager-markets},~\ref{fig:manager-users}.

\subsection{Проектирование пользовательского интерфейса}

Интерфейс спроектирован по принципам, сформулированным в разделе~1: фильтры отделены от области результатов, обратная связь при изменении параметров --- немедленная, предусмотрены светлая и тёмная темы.

\subsubsection{Структура основного экрана}

Экран дашборда делится на шапку и рабочую зону. В шапке --- выбор города, переключатель источника данных (для аналитика), элементы темы, имя пользователя и выход.

Боковая панель фильтров на десктопе закреплена слева; на узком экране открывается поверх контента. На рис.~\ref{fig:dashboard-filters} показаны слайдеры диапазонов, чипы комнатности и типа дома, переключатели исключения этажей и конструктор дополнительных условий.

\screenshot[0.78]{dashboard-filters}{dashboard-filters}{Панель фильтров дашборда с активными параметрами отбора}

Первый блок основной колонки --- \textbf{обзор рынка}: медианы цены, цены за м², площадей, число объектов, краткие текстовые инсайты по активным фильтрам. Ниже располагается \textbf{сетка графиков} --- гистограммы, столбчатые диаграммы по категориям, точечные диаграммы с линией тренда; у аналитика и администратора доступен конструктор пользовательских графиков. Общий вид этих блоков для роли \texttt{analyst} с серверным датасетом Москвы --- на рис.~\ref{fig:dashboard-main}.

\screenshot{dashboard-main}{dashboard-main}{Основной экран дашборда (роль analyst, город Москва): шапка, KPI и сетка графиков}

Далее по вертикали следует \textbf{карта} --- кластеры объявлений на базе JavaScript API Яндекс.Карт, балуны с ценой и ссылкой, ограничение числа отображаемых точек для производительности (рис.~\ref{fig:dashboard-map}).

\screenshot{dashboard-map}{dashboard-map}{Карта объектов с кластерами и открытым балуном объявления}

Блок \textbf{сравнения сегментов} выводит два независимых набора фильтров и KPI рядом --- например, новостройки и вторичка (рис.~\ref{fig:dashboard-segments}).

\screenshot{dashboard-segments}{dashboard-segments}{Сравнение двух сегментов рынка (два набора фильтров и KPI)}

Завершает экран \textbf{таблица объявлений} для ролей analyst и admin: сортировка, настраиваемые колонки, геокодирование строк без адреса, экспорт (рис.~\ref{fig:dashboard-table}).

\screenshot[0.95]{dashboard-table}{dashboard-table}{Таблица отфильтрованных объявлений с настраиваемыми колонками (роль analyst)}

\subsubsection{Типовые сценарии}

\textbf{Аналитик:} вход $\rightarrow$ выбор серверного снимка или загрузка CSV $\rightarrow$ настройка фильтров (рис.~\ref{fig:dashboard-filters}) $\rightarrow$ просмотр медиан и графиков (рис.~\ref{fig:dashboard-main}) $\rightarrow$ проверка расположения на карте (рис.~\ref{fig:dashboard-map}) $\rightarrow$ экспорт таблицы или графика (рис.~\ref{fig:dashboard-table}).

\textbf{Наблюдатель:} вход $\rightarrow$ просмотр серверного датасета без операций записи и без полной таблицы. Интерфейс ограничен KPI, графиками и картой (рис.~\ref{fig:observer-dashboard}).

\screenshot{observer-dashboard}{observer-dashboard}{Дашборд роли observer: KPI, графики и карта без таблицы и экспорта}

\textbf{Администратор:} вход $\rightarrow$ панель серверных данных --- upload, ingest, удаление (рис.~\ref{fig:admin-server}) $\rightarrow$ запуск стадии парсера и мониторинг прогресса (рис.~\ref{fig:admin-parser}) $\rightarrow$ при капче --- ручное продолжение в браузере парсера $\rightarrow$ ingest готового CSV.

\screenshot{admin-server}{admin-server}{Админ-панель управления серверным датасетом (upload, ingest, метаданные города)}

\screenshot{admin-parser}{admin-parser}{Админ-панель запуска парсера и отображения прогресса задачи}

\textbf{Руководитель:} вход в консоль manager $\rightarrow$ вкладка «Рынки» --- объём данных, даты обновления, диапазон цен по городам (рис.~\ref{fig:manager-markets}) $\rightarrow$ вкладка «Пользователи» --- создание учёток, смена ролей (рис.~\ref{fig:manager-users}).

\screenshot{manager-markets}{manager-markets}{Консоль руководителя: вкладка «Рынки» со сводкой по городам}

\screenshot[0.88]{manager-users}{manager-users}{Консоль руководителя: вкладка «Пользователи» (создание учётных записей и ролей)}

\subsection{Обоснование технологического стека}

Выбор средств разработки согласован с требованиями раздела~1 и архитектурой на рис.~\ref{fig:architecture}. Сводка приведена в таблице~\ref{tab:stack}.

\begin{table}[htbp]
\centering
\small
\setlength{\tabcolsep}{3pt}
\begin{tabular}{|p{0.14\textwidth}|p{0.19\textwidth}|p{0.61\textwidth}|}
\hline
\textbf{Уровень} & \textbf{Технология} & \textbf{Обоснование} \\
\hline
Парсер & Python, Selenium~4, pandas & Динамический SPA Яндекс.Недвижимость; ETL в табличном виде~\cite{selenium,mckinney2023} \\
\hline
API & Node.js, Express~5, better-sqlite3 & Лёгкий REST-сервис, JWT, синхронный доступ к SQLite без отдельного сервера БД \\
\hline
Клиент & React~19, TypeScript, Vite~6 & SPA с типизацией и быстрой пересборкой~\cite{react,typescript,vite,sultanov2022} \\
\hline
Графики & Recharts~3 & Декларативные React-компоненты, синхронная перерисовка при смене фильтров~\cite{recharts} \\
\hline
Таблицы/CSV & PapaParse, SheetJS & Разбор и экспорт CSV на клиенте без нагрузки на сервер~\cite{papaparse} \\
\hline
Стили & Tailwind CSS & Утилитарная вёрстка, адаптивность, темы~\cite{tailwind} \\
\hline
Карта & Яндекс.Карты JS API~2.1 & Кластеризация и балуны; согласованность с экосистемой источника данных \\
\hline
Среда & VS Code, WebStorm & Мультиязычный проект (Python + TypeScript), единый репозиторий \\
\hline
\end{tabular}
\caption{Технологический стек и мотивировка выбора}
\label{tab:stack}
\end{table}

\textbf{Selenium вместо статического HTML-парсинга} --- карточки и выдача на \texttt{realty.yandex.ru} формируются JavaScript; без управляемого браузера выборка неполна (ср.~табл.~1.1). \textbf{SQLite вместо PostgreSQL} --- достаточно для учебного развёртывания и демонстрации ролевой модели; объём объявлений вынесен в JSON. \textbf{Фильтрация на клиенте} --- после bootstrap сервер не участвует в каждом изменении слайдера, что соответствует паттерну SPA~\cite{nesterov2024,filimonova2024}.

\subsection{Выводы по второму разделу}

Во втором разделе спроектирован программный комплекс для анализа рынка недвижимости на данных Яндекс.Недвижимость:
\begin{enumerate}
    \item сформулированы функциональные и нефункциональные требования, включая поэтапный сбор, мультгородность и ролевую модель;
    \item выделены четыре логических компонента (клиент, API, парсер, общий слой правил) и описаны потоки данных между ними;
    \item спроектированы модель хранения (JSON-снимки + SQLite), контракт полей датасета и группы REST-эндпоинтов;
    \item определены экраны и сценарии для ролей observer, analyst, admin и manager; показано подключение API Яндекс.Карт и Геокодера;
    \item обоснован технологический стек с опорой на теоретический анализ первого раздела.
\end{enumerate}


\section{Реализация программного комплекса}

Во втором разделе описаны архитектура, модель данных и интерфейсы системы. Раздел посвящён программной реализации: Python-пайплайн сбора, серверное API на Node.js, общий слой правил в TypeScript и React-клиент. Каталоги \texttt{parser/}, \texttt{server/}, \texttt{shared/} и \texttt{src/} соответствуют архитектурным слоям системы.

\newcounter{codelisting}
\newcommand{\codelistingcaption}[1]{%
  \refstepcounter{codelisting}%
  \par\vspace{1.5\baselineskip}%
  \nopagebreak[4]%
  \begingroup\sloppy\raggedright\hbadness=10000\emergencystretch=3em
  \noindent\textbf{Листинг~3.\arabic{codelisting}.} #1\par
  \endgroup\vspace{0.8\baselineskip}%
}

% ---------------------------------------------------------------------------
\subsection{Общая структура репозитория и сборка}
% ---------------------------------------------------------------------------

\subsubsection{Каталоги и назначение}

Монорепозиторий \texttt{real-estate-analytics-dashboard} объединяет все компоненты аналитической системы. Корневой \texttt{package.json} описывает клиент (Vite~+~React) и скрипты запуска сервера; Python-парсер живёт в подкаталоге \texttt{parser/} со своим виртуальным окружением \texttt{parser/.venv} и файлом \texttt{requirements-parser.txt}.

\begin{itemize}
    \item \texttt{parser/} --- Selenium-парсер Яндекс.Недвижимость, ETL, QA, JSON-конфиги городов (\texttt{configs/}), артефакты сбора (\texttt{data/\{city\}/});
    \item \texttt{server/} --- Express~5, SQLite (\texttt{app.db}), хранение датасетов по городам, запуск задач парсера, обратное геокодирование;
    \item \texttt{shared/} --- типы, фильтры, профили городов, слияние строк; импортируется и сервером, и клиентом;
    \item \texttt{src/} --- SPA дашборда: фильтры, Recharts, карта Yandex Maps, админ-панели.
\end{itemize}

Такое разделение позволяет независимо обновлять сбор данных (Python) и визуализацию (TypeScript), сохраняя единые правила фильтрации в \texttt{shared/}.

\subsubsection{Скрипты npm и локальный запуск}

Сборка и разработка orchestrated через npm-скрипты корневого манифеста (листинг~3.25). Команда \texttt{npm run dev:all} параллельно поднимает API (\texttt{tsx server/index.ts}, порт~3001) и Vite-dev-server клиента (порт~5173). Отдельно доступны \texttt{npm run build} (production-бандл), \texttt{npm run check:parser-configs} (валидация JSON-конфигов парсера) и \texttt{npm run dev:api} только для backend.

Переменные окружения задаются в файле \texttt{.env} в корне репозитория: \texttt{JWT\_SECRET}, \texttt{YANDEX\_GEOCODER\_API\_KEY}, \texttt{VITE\_YANDEX\_MAPS\_API\_KEY}. Сервер при старте загружает их через \texttt{bootstrapEnv.ts}; без ключей карта и платный геокодер деградируют к публичным fallback-провайдерам (см.~листинг~3.14).

\subsubsection{Python-окружение парсера}

Парсер не требует глобальной установки пакетов: при первом запуске из админ-панели \texttt{parserRunner.ts} создаёт \texttt{parser/.venv} и выполняет \texttt{pip install -r requirements-parser.txt}. Альтернатива для ручной отладки --- \texttt{make moscow-links}, \texttt{make moscow-details} и аналоги в \texttt{parser/Makefile}, которые активируют venv и вызывают модули \texttt{pipeline.run\_*}.

Клиентская часть собирается Vite с \texttt{@vitejs/plugin-react}. \texttt{tsconfig.json} задаёт paths \texttt{@/*}; серверный \texttt{tsx} резолвит \texttt{../shared/} без отдельной компиляции общего слоя.

Конфигурация города --- JSON (\srce{parser/configs/config.moscow.json} и аналоги): URL сегментов, координаты, пути к CSV, Selenium, лимиты. \texttt{PipelineConfig} (листинг~3.1) загружает файл и собирает URL для стадии \emph{links}.

\codelistingcaption{Загрузка конфигурации пайплайна (\srce{parser/pipeline/config.py})}
\begin{singleminted}{python}
@dataclass
class PipelineConfig:
    base_url: str
    search_segments: list[str] | None
    headless: bool
    target_listings: int | None
    links_file: str
    details_file: str
    processed_file: str
    resume_details: bool
    resume_after_captcha: bool
    runtime_status_file: str
    runtime_control_file: str
    city_center_lat: float
    city_center_lng: float
    request_delays_ms: dict[str, dict[str, int]]

    @staticmethod
    def from_file(
        path: str | Path,
        *,
        preset_override: str | None = None,
    ) -> "PipelineConfig":
        cfg_path = Path(path)
        raw = json.loads(cfg_path.read_text(encoding="utf-8"))
        raw = apply_preset(raw)
        # ... загрузка search_segments и search_segments_file ...
        return PipelineConfig(
            base_url=raw.get("base_url", "https://realty.yandex.ru/..."),
            headless=bool(raw.get("headless", True)),
            resume_details=bool(raw.get("resume_details", True)),
            city_center_lat=float(raw.get("city_center_lat", 51.533103)),
            # ...
        )

    @property
    def search_urls(self) -> list[str]:
        urls: list[str] = []
        if self.base_urls:
            urls.extend(self.base_urls)
        else:
            urls.append(self.base_url)
        if self.search_segments:
            urls.extend(self.search_segments)
        return _dedupe_urls(urls)
\end{singleminted}
\label{lst:pipeline-config}

Фрагмент демонстрирует централизованное чтение JSON: все стадии получают один объект конфигурации, а \texttt{search\_urls} объединяет базовые URL и сегменты (для Москвы --- \srce{moscow_search_segments.json}).

Для Москвы конфигурация включает двенадцать базовых URL плюс дополнительные сегменты из JSON. Параметр \texttt{max\_pages=150} ограничивает глубину пагинации; \texttt{target\_listings} останавливает сбор после нужного объёма.

Артефакты сбора складываются в \srce{parser/data/moscow/} или \srce{parser/data/saratov/}: список ссылок, details CSV, \srce{processed_apartment_data.csv}, статистика парсинга.

% ---------------------------------------------------------------------------
\subsection{Реализация модуля парсера}
% ---------------------------------------------------------------------------

\subsubsection{Архитектура Python-пайплайна}

Парсер --- цепочка CLI-модулей в \texttt{pipeline/}. Оркестратор \srce{run_pipeline.py} (листинг~3.6) вызывает \texttt{run\_links}, \texttt{run\_details}, \texttt{etl}, \texttt{qa\_report} через \texttt{subprocess}, обновляя \srce{.runtime_status.json}.

Браузер обслуживают \srce{yandex_realty_advanced_parser.py} (список) и \srce{yandex_realty_details_parser.py} (карточки) на Selenium~4 с \texttt{webdriver-manager}.

Details-парсер извлекает JSON-LD и поля «О доме», «Удобства», «Двор», нормализует число комнат и записывает геолокацию. Особое внимание уделено устойчивости записи и интеграции с сервером.

\subsubsection{Стадия links: сбор URL карточек}

Скрипт \texttt{pipeline/run\_links.py} (листинг~3.2) для каждого поискового сегмента из конфигурации создаёт экземпляр \texttt{YandexRealtyParser}, прокручивает выдачу, переключает страницы пагинации и дописывает уникальные ссылки в \texttt{yandex\_realty\_links.txt}. Между scroll, navigation и pagination действуют случайные паузы из блока \texttt{request\_delays\_ms.links}, что снижает частоту срабатывания антибот-защиты.

При достижении \texttt{target\_listings} цикл по сегментам прерывается; если файл уже содержит достаточно ссылок, стадия завершается без запуска браузера. Callback \texttt{on\_status\_update} записывает в \texttt{.runtime\_status.json} поля \texttt{stage}, \texttt{done}, \texttt{total}, \texttt{statusMessage} --- их потребляет серверная задача парсера.

\codelistingcaption{Точка входа стадии links (\srce{parser/pipeline/run_links.py})}
\begin{singleminted}{python}
def main() -> None:
    cfg = PipelineConfig.from_file(args.config)
    links_path = cfg.resolve(cfg.links_file)
    runtime_status_path = cfg.resolve(cfg.runtime_status_file)
    runtime_control_path = cfg.resolve(cfg.runtime_control_file)
    clear_runtime_control(runtime_control_path)

    p = links_parser.YandexRealtyParser(
        headless=cfg.headless,
        max_pages=cfg.max_pages,
        target_listings=target_listings,
        request_delays_ms=cfg.request_delays_ms.get("links"),
    )
    p.links_file = str(links_path)
    p.should_resume = lambda: consume_runtime_resume(runtime_control_path)

    def on_status(payload: dict[str, Any]) -> None:
        write_runtime_status(runtime_status_path, {
            "stage": "links",
            "statusMessage": payload.get("statusMessage") or "Сбор ссылок",
            **payload,
        })
    p.on_status_update = on_status

    for idx, url in enumerate(cfg.search_urls, start=1):
        if target_listings and current_count >= target_listings:
            break
        links_parser.BASE_URL = url
        write_runtime_status(runtime_status_path, {
            "stage": "links", "status": "running",
            "step": f"Сегмент {idx}/{len(urls)}",
        })
        p.run()
\end{singleminted}
\label{lst:run-links}

\subsubsection{Стадия details: извлечение атрибутов карточки}

Модуль \texttt{yandex\_realty\_details\_parser.py} обходит список ссылок и из DOM извлекает цену, площади, адрес, координаты, параметры дома, удобства. Класс \texttt{DetailsStore} (\texttt{pipeline/details\_store.py}, листинг~3.3) обеспечивает \emph{безопасную} запись: промежуточный буфер \texttt{.tmp\_details\_output.csv}, атомарный merge в основной CSV, JSONL-журнал построчных результатов и резервные копии с суффиксом \texttt{.latest}. При \texttt{resume\_details=true} повторно обрабатываются только URL, отсутствующие в store --- это критично при сбоях сети или капче на тысячах карточек.

\codelistingcaption{Безопасная запись details CSV (\srce{parser/pipeline/details_store.py})}
\begin{singleminted}{python}
def write_rows(path: Path, rows: List[Dict[str, str]]) -> None:
    path.parent.mkdir(parents=True, exist_ok=True)
    if not rows:
        raise RuntimeError(f"Отказ записи пустого CSV: {path}")

    expected = unique_link_count(rows)
    fieldnames = details_fieldnames(rows)
    part = path.with_suffix(path.suffix + ".part")

    with part.open("w", encoding="utf-8", newline="") as f:
        writer = csv.DictWriter(f, fieldnames=fieldnames, extrasaction="ignore")
        writer.writeheader()
        writer.writerows(rows)

    written = count_links_in_file(part)
    if written < expected:
        part.unlink(missing_ok=True)
        raise RuntimeError(
            f"Запись оборвалась: ожидали {expected} строк, в part {written}"
        )
    part.replace(path)
    # ... копия в .latest и jsonl ...
\end{singleminted}
\label{lst:details-store}

Проверка \texttt{written < expected} защищает от «обнуления» файла при обрыве записи --- типичная проблема длительных сессий парсинга.

\subsubsection{Обнаружение SmartCaptcha и возобновление}

Яндекс периодически показывает SmartCaptcha. Парсер не пытается решать её автоматически: функция \texttt{should\_treat\_as\_captcha\_page} (листинг~3.4) анализирует URL, заголовок страницы и видимые DOM-маркеры, но \emph{игнорирует} скрытые iframe, если карточка объявления уже загружена --- это снижает ложные срабатывания (покрыто unit-тестами в \texttt{test\_captcha\_detection.py}).

При подтверждённой капче метод \texttt{handle\_captcha} (листинг~3.5) переключает Chrome из headless в оконный режим, пишет статус \texttt{captcha\_required} в runtime JSON и ждёт действия пользователя. В неинтерактивном режиме (запуск из админ-панели) цикл опрашивает \texttt{.runtime\_control.json}: сервер по \texttt{POST /api/admin/parser/jobs/:id/resume} вызывает \texttt{request\_runtime\_resume}, Python --- \texttt{consume\_runtime\_resume} (листинг~3.7).

\codelistingcaption{Детектор страницы капчи (\srce{parser/yandex_realty_details_parser.py})}
\begin{singleminted}{python}
def should_treat_as_captcha_page(
    *,
    url: str = "",
    page_title: str = "",
    offer_card_present: bool = False,
    visible_captcha_marker: bool = False,
) -> bool:
    if offer_card_present and not url_indicates_captcha_page(url):
        return False
    if url_indicates_captcha_page(url):
        return True
    if page_title_indicates_captcha(page_title):
        return True
    return visible_captcha_marker
\end{singleminted}
\label{lst:captcha-detect}

\codelistingcaption{Обработка капчи и resume-сигнала (\srce{parser/yandex_realty_details_parser.py})}
\begin{singleminted}{python}
def handle_captcha(self, url: str = "", link_index: int | None = None):
    print("PARSER_PROGRESS_STATUS:captcha_required")
    self.emit_status({"status": "captcha_required", "step": "Капча"})
    if self.headless:
        self.headless = False
        self.recreate_driver("капча: нужно ручное решение")
        if url:
            self.driver.get(url)

    if not self.interactive:
        started = time.time()
        while True:
            if self.should_resume and self.should_resume():
                print("Сигнал продолжения из UI.")
            if not self.check_for_captcha():
                self.emit_status({"status": "running", "step": "Капча решена"})
                return True
            if time.time() - started >= self.captcha_auto_wait_seconds:
                self.emit_status({"status": "waiting_user", ...})
            time.sleep(3)
\end{singleminted}
\label{lst:captcha-handle}

\subsubsection{Стадия ETL: нормализация и производные признаки}

Модуль \texttt{pipeline/etl.py} преобразует \texttt{yandex\_realty\_details.csv} в \texttt{processed\_apartment\_data.csv}, согласованный с \texttt{schema\_contract.json}. На входе --- «сырые» строки details; на выходе --- таблица, готовая для ingest на сервер и загрузки в дашборд.

Ключевые шаги реализованы в функциях:
\begin{itemize}
    \item \texttt{to\_float}, \texttt{to\_int}, \texttt{parse\_rooms} --- приведение числовых полей и извлечение числа комнат из заголовка;
    \item \texttt{sanitize\_address} --- удаление служебных фраз, координат вместо адреса, «загрязнённого» адреса офиса Яндекса на Садовнической (листинг~3.26);
    \item \texttt{haversine\_km} --- расстояние до центра города из конфигурации;
    \item \texttt{build\_row} --- расчёт цены за м², возраста дома, относительного этажа, бинарных признаков удобств;
    \item \texttt{apply\_one\_hot} --- кодирование категорий через \texttt{categories.py} (санузел, окна, ремонт, тип дома).
\end{itemize}

Функция \texttt{build\_row} (листинг~3.8) связывает сырые и итоговые поля. Модуль \texttt{categories.py} задаёт \texttt{CATEGORY\_CLASSIFIERS}: для каждого текстового поля указана функция классификации; \texttt{apply\_one\_hot} (листинг~3.25) разворачивает категории в булевы колонки с префиксами \texttt{тип\_дома\_}, \texttt{ремонт\_} и др.

После построения всех строк выполняется дедупликация по канонической ссылке или ID объявления: если одна карточка попала в details дважды (например, после resume), сохраняется последняя версия строки.

\codelistingcaption{Очистка текстового адреса в ETL (\srce{parser/pipeline/etl.py})}
\begin{singleminted}{python}
def sanitize_address(raw) -> str:
    if pd.isna(raw):
        return ""
    normalized = str(raw).strip().replace("\xa0", " ")
    normalized = re.sub(
        r"^(?:расположение|адрес)\s*[:\-]\s*", "", normalized, flags=re.I
    ).strip(" \"'«»")
    if NUMERIC_ONLY_RE.fullmatch(normalized):
        return ""
    if COORD_PAIR_RE.fullmatch(normalized):
        return ""
    if is_polluted_parser_address(normalized):
        return ""
    if not ADDRESS_HAS_LETTER_RE.search(normalized):
        return ""
    return normalized
\end{singleminted}
\label{lst:sanitize-address}

\codelistingcaption{Построение строки processed CSV (\srce{parser/pipeline/etl.py})}
\begin{singleminted}{python}
def build_row(row, current_year, center_lat, center_lng, output_columns):
    listing_id = extract_listing_id(row)
    listing_link = normalize_listing_link(row.get("Ссылка"), listing_id)
    price = to_float(row.get("Цена"))
    area = to_float(row.get("Общая площадь"))
    lat, lng = parse_geo(row.get("Геолокация"))
    distance = (
        haversine_km(lat, lng, center_lat, center_lng)
        if lat is not None and lng is not None else None
    )
    price_per_m2 = (price / area) if price and area else None
    floor_rel, first_floor, last_floor = parse_floor(
        row.get("Этаж"), row.get("Этажей в доме")
    )
    location = sanitize_address(row.get("Расположение"))
    out = {
        "ID": listing_id,
        "Расположение": location,
        "Ссылка": listing_link,
        "Цена": price,
        "Широта": lat, "Долгота": lng,
        "Расстояние до центра (км)": distance,
        "Цена за м²": price_per_m2,
        "Этаж_относительный": floor_rel,
        "Первый_этаж": first_floor,
        "Последний_этаж": last_floor,
    }
    apply_one_hot(out, row, output_columns)
    return out
\end{singleminted}
\label{lst:etl-build-row}

\subsubsection{Контроль качества (QA)}

\texttt{pipeline/qa\_report.py} формирует отчёт о качестве итогового CSV: сверка с контрактом схемы (\texttt{validate\_schema}, листинг~3.22), доля пропусков в обязательных колонках, число дубликатов ссылок, заполненность координат. Контракт \texttt{schema\_contract.json} перечисляет \texttt{output\_columns} и \texttt{required\_core}; ETL и QA используют один и тот же файл, поэтому рассогласование схемы обнаруживается до попадания данных в дашборд.

Результаты сохраняются в JSON и Markdown в каталоге данных города; при критических ошибках администратор видит их до ingest на сервер. Отчёт включает топ колонок по доле пропусков, счётчик строк с валидными широтой и долготой, предупреждение при заполненности ID ниже 95\%.

\subsubsection{Оркестрация и файлы статуса}

\codelistingcaption{Pipeline и обновление статуса (\srce{parser/pipeline/run_pipeline.py})}
\begin{singleminted}{python}
def run_step(module_name: str, config_path: str, *, target_listings=None) -> None:
    cmd = [sys.executable, "-m", module_name, "--config", config_path]
    if target_listings and module_name in {"pipeline.run_links", "pipeline.run_details"}:
        cmd.extend(["--target-listings", str(target_listings)])
    subprocess.run(cmd, check=True, env=os.environ.copy())

def main() -> None:
    cfg = PipelineConfig.from_file(config_abs)
    if not args.skip_links:
        write_runtime_status(runtime_status_path, {"stage": "links", "status": "running"})
        run_step("pipeline.run_links", config_abs, target_listings=args.target_listings)
    if not args.skip_details:
        run_step("pipeline.run_details", config_abs, ...)
    if not args.skip_etl:
        run_step("pipeline.etl", config_abs)
    if not args.skip_qa:
        run_step("pipeline.qa_report", config_abs)
\end{singleminted}
\label{lst:run-pipeline}

\codelistingcaption{Протокол resume через JSON (\srce{parser/pipeline/runtime_status.py})}
\begin{singleminted}{python}
def request_runtime_resume(path: Path) -> None:
    current = _read_json(path)
    current["resume_requested_at"] = now_iso()
    path.write_text(json.dumps(current, ensure_ascii=False, indent=2), encoding="utf-8")

def consume_runtime_resume(path: Path) -> bool:
    current = _read_json(path)
    stamp = current.get("resume_requested_at")
    if not stamp:
        return False
    current["resume_requested_at"] = None
    path.write_text(json.dumps(current, ensure_ascii=False, indent=2), encoding="utf-8")
    return True
\end{singleminted}
\label{lst:runtime-resume}

% ---------------------------------------------------------------------------
\subsection{Реализация серверной части}
% ---------------------------------------------------------------------------

\subsubsection{Express-приложение и инициализация}

Точка входа --- \texttt{server/index.ts}. При старте создаётся каталог \texttt{server/data/}, открывается SQLite с режимом WAL, выполняется миграция таблицы пользователей и seed-учётки (\texttt{admin}, \texttt{analyst}, \texttt{observer}, \texttt{manager}). Инициализируются \texttt{CityDatasetManager}, \texttt{ParserJobService} и \texttt{ReverseGeocoderService}.

Маршруты сгруппированы по префиксам: \texttt{/api/auth/*}, \texttt{/api/dataset/*}, \texttt{/api/listings}, \texttt{/api/geocode/*}, \texttt{/api/admin/*}, \texttt{/api/manager/*}. Параметр \texttt{city=moscow|saratov} обязателен для операций с датасетом.

Класс \texttt{CityDatasetManager} --- фасад над экземплярами \texttt{CityDatasetStore} (по одному на город). Метод \texttt{listCities()} возвращает метку, наличие данных и число строк. SQLite хранит пользователей, задачи парсера и кэш геокодера; объявления --- в JSON.

\subsubsection{Аутентификация и авторизация}

Аутентификация построена на JWT. Эндпоинт \texttt{POST /api/auth/login} (листинг~3.9) сравнивает пароль через \texttt{bcryptjs}, выдаёт токен на 12~часов. Middleware \texttt{authMiddleware} извлекает Bearer-токен; \texttt{requireRoles([...])} ограничивает admin-, analyst- и manager-эндпоинты.

\codelistingcaption{JWT middleware и вход (\srce{server/index.ts})}
\begin{singleminted}{typescript}
function authMiddleware(req, res, next): void {
    const auth = req.headers.authorization;
    if (!auth?.startsWith('Bearer ')) {
        res.status(401).json({ error: 'Unauthorized' });
        return;
    }
    try {
        const token = auth.slice(7);
        const user = jwt.verify(token, JWT_SECRET) as AuthUser;
        (req as AuthRequest).user = user;
        next();
    } catch {
        res.status(401).json({ error: 'Invalid token' });
    }
}

app.post('/api/auth/login', (req, res) => {
    const row = db.prepare('SELECT ... FROM users WHERE username = ?').get(username);
    if (!row || !bcrypt.compareSync(password, row.password_hash)) {
        res.status(401).json({ error: 'Неверный логин или пароль' });
        return;
    }
    const token = signToken({
        id: row.id, username: row.username, role: row.role,
    });
    res.json({
        token,
        user: { id: row.id, username: row.username, role: row.role },
    });
});
\end{singleminted}
\label{lst:auth}

\subsubsection{Хранение и ingest датасетов}

Класс \texttt{CityDatasetStore} в \texttt{server/datasetStore.ts} управляет JSON-файлами \texttt{server/data/cities/\{cityId\}/dataset.json} и метаданными загрузок. Метод \texttt{ingestCsv} (листинг~3.10):
\begin{enumerate}
    \item разбирает текст CSV через PapaParse;
    \item отфильтровывает строки с координатами за пределами радиуса города (\texttt{getCityDistanceLimitKm});
    \item в режиме \texttt{replace} полностью заменяет датасет; в \texttt{append} --- сливает строки с одинаковым ID через \texttt{mergeTwoDatasetRows};
    \item присваивает \texttt{batchId} для возможности отката партии.
\end{enumerate}

Эндпоинт \texttt{POST /api/admin/dataset/ingest-pipeline} импортирует готовый \texttt{processed\_apartment\_data.csv} из каталога парсера без ручной загрузки файла --- связующее звено между Python и Node.

\codelistingcaption{Ingest CSV на сервере (\srce{server/datasetStore.ts})}
\begin{singleminted}{typescript}
ingestCsv(csvText: string, fileName: string, mode: UploadMode) {
    const parsed = parseCsvText(csvText);
    const filteredParsed = filterRowsTooFarFromCityCenter(
        parsed.rows, this.cityProfile
    );
    const batchId = randomUUID();

    if (mode === 'replace') {
        const mergedRows = applyRowMerge({
            rows: filteredParsed.rows,
            summary: buildSummary(filteredParsed.rows),
        }).state.rows;
        const rows = tagRows(validatedRows.rows, batchId);
        this.state = { rows, summary: buildSummary(rows) };
        this.meta.uploads = [{ batchId, at, fileName, added: rows.length, ... }];
    } else {
        // append: mergeTwoDatasetRows по offer id
    }
    this.persist();
    return { totalRows, added, skippedDuplicates, tooFarFiltered, ... };
}
\end{singleminted}
\label{lst:ingest-csv}

\subsubsection{Запуск парсера из Node.js}

Модуль \texttt{parserRunner.ts} связывает админ-панель с Python. Функция \texttt{spawnParserJob} (листинг~3.11) выполняет preflight-валидацию конфига, разрешает интерпретатор Python (venv или bootstrap), формирует аргументы \texttt{-m pipeline.run\_*} и запускает \texttt{spawn} с перенаправлением stdout/stderr в лог задачи.

Соответствие стадий модулям:
\begin{itemize}
    \item \texttt{links} $\rightarrow$ \texttt{pipeline.run\_links};
    \item \texttt{details} $\rightarrow$ \texttt{pipeline.run\_details};
    \item \texttt{flush} $\rightarrow$ \texttt{run\_details --flush-only};
    \item \texttt{etl} $\rightarrow$ \texttt{pipeline.etl};
    \item \texttt{pipeline} $\rightarrow$ \texttt{pipeline.run\_pipeline}.
\end{itemize}

Состояние job хранится в \texttt{parser\_jobs}. UI опрашивает \srce{GET /api/admin/parser/jobs/:id} и читает прогресс из \srce{.runtime_status.json}.

Сервис \texttt{ParserJobService} регистрирует запуск: UUID, город, стадию, путь к логу, PID. \texttt{resumeCaptcha} пишет в \srce{.runtime_control.json} метку \texttt{resume\_requested\_at}. Отмена посылает SIGTERM дочернему процессу.

Перед первым запуском для города выполняется \texttt{validateParserConfig}: проверяется существование JSON, соответствие URL маркерам города (\texttt{/moskva/}, \texttt{/saratov/}), число сегментов поиска и пути к артефактам. Ошибки preflight возвращаются в UI до spawn, что экономит время на «пустых» задачах с неверным конфигом.

\codelistingcaption{Запуск Python-процесса парсера (\srce{server/parserRunner.ts})}
\begin{singleminted}{typescript}
export function spawnParserJob(cityId, stage, logPath, options?) {
    const preflight = validateParserConfig(cityId);
    if (!preflight.ok) {
        return { error: preflight.errors.join('\n') };
    }

    const python = resolveParserPython(stage);
    if (!python.ok) return { error: python.error };

    const moduleName = STAGE_MODULE[stage]; // pipeline.run_links, ...
    const args = ['-m', moduleName, '--config', configPath];
    if (stage === 'flush') args.push('--flush-only');
    if (targetListings != null) args.push('--target-listings', String(targetListings));

    const child = spawn(python.pythonBin, args, {
        cwd: PARSER_DIR,
        env: { ...process.env, PYTHONPATH: PARSER_DIR },
        stdio: ['ignore', 'pipe', 'pipe'],
    });
    child.stdout?.pipe(logStream);
    child.stderr?.pipe(logStream);
    return { child, configPath };
}
\end{singleminted}
\label{lst:parser-runner}

\subsubsection{Обратное геокодирование}

Сервис \texttt{ReverseGeocoderService} (\texttt{reverseGeocoder.ts}) по паре (широта, долгота) возвращает текстовый адрес. Цепочка провайдеров (листинг~3.12): in-memory кэш $\rightarrow$ SQLite $\rightarrow$ Yandex Geocoder API (ключ \texttt{YANDEX\_GEOCODER\_API\_KEY}) $\rightarrow$ публичный Yandex $\rightarrow$ Nominatim (OpenStreetMap). Ключ кэша формируется в \texttt{shared/reverseGeocodeKey.ts} с округлением координат до 5 знаков; TTL по умолчанию --- 14 суток.

Ключи API получены в кабинете разработчика Яндекса (рис.~\ref{fig:yandex-api}, раздел~2): JavaScript API --- для карты на клиенте, Geocoder API --- для серверного кэша адресов.

\codelistingcaption{Цепочка reverse geocode (\srce{server/reverseGeocoder.ts})}
\begin{singleminted}{typescript}
private async reverseWithCache(lat, lng, cacheKey) {
    const dbRow = this.lookupCacheRow(cacheKey, lat, lng);
    if (dbRow?.address && dbRow.expires_at > now) {
        return { cacheKey, address: dbRow.address, source: 'cache', provider: 'cache' };
    }
    if (this.apiKey) {
        const yandexWithKey = await this.fetchYandexAddress(lat, lng, this.apiKey);
        if (yandexWithKey.address) {
            this.saveCache(cacheKey, lat, lng, yandexWithKey.address);
            return { ..., provider: 'yandex-key' };
        }
    }
    const yandexWithoutKey = await this.fetchYandexAddress(lat, lng, null);
    if (yandexWithoutKey.address) { ... provider: 'yandex-public' }

    const nominatim = await this.fetchNominatimAddress(lat, lng);
    if (nominatim.address) { ... provider: 'nominatim-public' }
    return { cacheKey, address: null, source: 'fallback', ... };
}
\end{singleminted}
\label{lst:reverse-geocode}

% ---------------------------------------------------------------------------
\subsection{Общий слой shared}
% ---------------------------------------------------------------------------

Каталог \texttt{shared/} --- единый источник правил для сервера и клиента. Дублирование логики фильтрации в двух языках привело бы к расхождению KPI на карте и в таблице; поэтому \texttt{rowPassesFiltersShared} вызывается и в \texttt{GET /api/listings}, и в \texttt{App.tsx}.

\subsubsection{Профили городов}

Файл \texttt{cities.ts} (листинг~3.13) задаёт идентификаторы \texttt{moscow} и \texttt{saratov}, координаты центра карты, лимит расстояния для очистки данных (120 и 100~км) и путь к JSON-конфигу парсера. Сервер и клиент импортируют одни и те же константы.

\codelistingcaption{Профили городов (\srce{shared/cities.ts})}
\begin{singleminted}{typescript}
export const CITY_PROFILES: CityProfile[] = [
    {
        id: 'moscow',
        label: 'Москва',
        centerLat: 55.7558,
        centerLng: 37.6176,
        distanceMaxKm: 120,
        parserConfig: 'parser/configs/config.moscow.json',
    },
    {
        id: 'saratov',
        label: 'Саратов',
        centerLat: 51.533103,
        centerLng: 46.034266,
        distanceMaxKm: 100,
        parserConfig: 'parser/configs/config.saratov.json',
    },
];

export function getCityDistanceLimitKm(cityId?: string | null): number | null {
    return getCityProfile(cityId)?.distanceMaxKm ?? null;
}
\end{singleminted}
\label{lst:cities}

\subsubsection{Единая фильтрация строк}

Функция \texttt{rowPassesFiltersShared} (листинг~3.14) --- центральная точка отбора. Диапазон считается «активным» только если отличается от baseline summary (\texttt{isRangeFilterActive}): слайдер на полном диапазоне не отсекает данные. Проверяются цена, площадь, комнаты, год постройки, расстояние, тип дома, исключение первого/последнего этажа, радиус от центра (haversine при отсутствии колонки расстояния) и произвольные \texttt{additionalFilters}.

\codelistingcaption{Общая функция фильтрации (\srce{shared/filters.ts})}
\begin{singleminted}{typescript}
export function rowPassesFiltersShared(row, filters, summary): boolean {
    const c = summary.coreColumnMap;
    const priceActive = isRangeFilterActive(filters.price, summary.price, 0.5);
    const areaActive = isRangeFilterActive(filters.area, summary.area, 0.1);

    if (c.price && priceActive) {
        const v = row[c.price];
        if (!isFiniteNumber(v)) return false;
        if (v < filters.price.min || v > filters.price.max) return false;
    }
    if (c.rooms && filters.rooms.length > 0) {
        const rounded = parseRoomCount(row[c.rooms]);
        if (rounded == null || !filters.rooms.includes(rounded)) return false;
    }
    if (filters.radiusKm != null && c.distanceKm) {
        const v = asNumber(row[c.distanceKm]);
        if (!Number.isFinite(v) || v > filters.radiusKm) return false;
    }
    if (filters.excludeFirstFloor && c.firstFloor && asBool(row[c.firstFloor]))
        return false;
    // ... additionalFilters, houseTypes ...
    return true;
}
\end{singleminted}
\label{lst:filters}

\subsubsection{Слияние строк датасета}

При append или ingest строки с одним ID объявления объединяются модулем \texttt{mergeDatasetRows.ts} (листинг~3.15). Пустые поля заполняются из другой записи; HTTP-ссылка имеет приоритет. Типичный сценарий: одна выгрузка содержит координаты без ссылки, другая --- ссылку без координат; после merge строка пригодна и для карты, и для перехода на карточку.

\codelistingcaption{Слияние строк по offer id (\srce{shared/mergeDatasetRows.ts})}
\begin{singleminted}{typescript}
export function mergeDatasetRows(rows: DataRow[]) {
    const normalized = rows.map(normalizeDatasetRow);
    const groups = new Map<string, DataRow[]>();

    for (const row of normalized) {
        const ids = offerIdsFromRow(row);
        if (!ids.length) { noId.push(row); continue; }
        const key = ids[0];
        groups.get(key)?.push(row) ?? groups.set(key, [row]);
    }

    for (const group of groups.values()) {
        if (group.length === 1) { merged.push(group[0]); continue; }
        let acc = group[0];
        for (let i = 1; i < group.length; i++) {
            acc = mergeRowFields(acc, group[i]);
        }
        merged.push(acc);
    }
    return { rows: merged, stats: { inputRows, outputRows, mergedGroups, ... } };
}
\end{singleminted}
\label{lst:merge-rows}

% ---------------------------------------------------------------------------
\subsection{Реализация клиентской части}
% ---------------------------------------------------------------------------

\subsubsection{Маршрутизация экранов без react-router}

Корневой компонент \texttt{App.tsx} не использует react-router: экран выбирается по роли пользователя и фазе загрузки \texttt{loadPhase}. После входа manager видит \texttt{ManagerConsole}; admin при пустом городе --- \texttt{AdminSetupView}; analyst и observer --- основной дашборд после успешного \texttt{bootstrap}. Состояние \texttt{LoadPhase} явно кодирует ветвления: \texttt{pick-source}, \texttt{need-personal}, \texttt{server-empty}, \texttt{ready}.

Тип \texttt{DataSourceMode} различает источник данных \texttt{server} и \texttt{personal}. Аналитик может переключаться между серверной выгрузкой города и локальным CSV; выбор сохраняется в \texttt{localStorage}. Роль \texttt{observer} принудительно использует сервер; \texttt{manager} не попадает на дашборд вовсе. Переключатель города \texttt{CitySwitcher} вызывает \texttt{loadServerDataset} с новым \texttt{cityId} без перезагрузки страницы.

Функция \texttt{applyDataset} (листинг~3.16) --- центральный pipeline клиента: \texttt{mergeDatasetRows} $\rightarrow$ \texttt{applyListingQualityFilter} $\rightarrow$ пересчёт summary $\rightarrow$ восстановление фильтров и графиков из URL/localStorage. Фильтр качества отбрасывает строки без цены, без площади или с явно некорректными координатами --- последний барьер перед визуализацией, дублирующий часть проверок QA на стороне Python. Обновление обёрнуто в \texttt{startTransition}, чтобы UI оставался отзывчивым.

\codelistingcaption{Применение загруженного датасета (\srce{src/App.tsx})}
\begin{singleminted}{typescript}
const applyDataset = useCallback((data, summary, fileKey) => {
    const { rows: mergedRows } = mergeDatasetRows(data);
    const mergedSummary = refreshSummaryFromData(mergedRows, summary);
    const { rows: qualityRows } = applyListingQualityFilter(
        mergedRows, mergedSummary
    );
    const refreshed = refreshSummaryFromData(qualityRows, mergedSummary);
    const chosen = resolveDashboardStateForSummary(refreshed);
    startTransition(() => {
        setAllData(qualityRows);
        setDataSummary(refreshed);
        setFilters(chosen.filters);
        setBaselineFilters(createDefaultFiltersFromSummary(refreshed));
        setUserCharts(chosen.userCharts);
        setLoadPhase('ready');
    });
}, [startTransition, caps]);
\end{singleminted}
\label{lst:apply-dataset}

\subsubsection{Импорт CSV на клиенте}

Модуль \texttt{src/utils/csvImport.ts} (листинг~3.17) разбирает файл через PapaParse, обрабатывает строки чанками по 400 с вызовом \texttt{yieldToMain()} (\texttt{setTimeout(0)}) между чанками --- main thread не блокируется при 10~000+ строк. Лимиты совпадают с сервером: 20~МБ, 50~000 строк. Fingerprint файла (имя, размер, дата изменения) кэшируется в \texttt{localStorage} через модуль \texttt{cache.ts}: повторное открытие того же CSV мгновенно восстанавливает таблицу без повторного PapaParse. При ошибке формата пользователь получает сообщение на русском языке --- это согласовано с серверными кодами 413 и текстами ошибок ingest.

\codelistingcaption{Чанковый разбор CSV (\srce{src/utils/csvImport.ts})}
\begin{singleminted}{typescript}
const ROW_PROCESS_CHUNK = 400;

async function processRawRows(rawData, columnOrder, houseTypeCols) {
    const processed: DataRow[] = [];
    for (let i = 0; i < rawData.length; i += ROW_PROCESS_CHUNK) {
        const end = Math.min(i + ROW_PROCESS_CHUNK, rawData.length);
        for (let j = i; j < end; j += 1) {
            const out: DataRow = {};
            for (const k of columnOrder) {
                out[k] = normalizeCell(rawData[j][k]);
            }
            out.houseType = collectHouseType(rawData[j], houseTypeCols);
            if (rowHasAnyValue(out)) processed.push(out);
        }
        if (end < rawData.length) await yieldToMain();
    }
    return processed;
}
\end{singleminted}
\label{lst:csv-import}

\subsubsection{Фильтры и отложенное обновление UI}

Состояние фильтров хранится в \texttt{App.tsx} и передаётся в \texttt{FilterPanel}. Отфильтрованная выборка вычисляется в \texttt{useMemo}; для тяжёлых виджетов применяется \texttt{useDeferredValue(filters)} (листинг~3.18) --- KPI и графики обновляются с небольшой задержкой, поля ввода остаются отзывчивыми. Компонент \texttt{FilterPanel} предоставляет слайдеры диапазонов, чипы комнат и типов домов, переключатели «Искл.~1 этаж» / «Искл.~последний», а также конструктор дополнительных условий по произвольным колонкам (внешний вид --- рис.~\ref{fig:dashboard-filters}).

\codelistingcaption{Deferred-фильтрация и подсчёт (\srce{src/App.tsx})}
\begin{singleminted}{typescript}
const deferredFilters = useDeferredValue(filters);
const deferredUserCharts = useDeferredValue(userCharts);

const filteredData = useMemo(() => {
    if (!allData || !deferredFilters || !dataSummary) return [];
    const out: DataRow[] = [];
    for (const row of allData) {
        if (rowPassesFilters(row, deferredFilters, dataSummary)) out.push(row);
    }
    return out;
}, [allData, deferredFilters, dataSummary]);

const deferredFilteredData = useDeferredValue(filteredData);
const isChartUpdating =
    isPending || deferredFilteredData !== filteredData;
\end{singleminted}
\label{lst:deferred-filters}

\subsubsection{Динамическая сетка графиков}

Компонент \texttt{DynamicChartGrid} строит гистограммы, scatter и bar chart по конфигурации \texttt{UserChartDefinition[]}. Перед отрисовкой для каждого типа графика проверяется «допустимость» столбца: модуль \texttt{chartColumnRules.ts} запрещает, например, scatter по категориальным one-hot полям или гистограмму по бинарным признакам \texttt{*\_есть}. Это предотвращает бессмысленные визуализации при конструкторе графиков.

Sidebar \texttt{ChartSidebar} (analyst/admin) добавляет графики, меняет порядок drag-and-drop, экспортирует PNG и PDF. Состояние сериализуется в shareable URL вместе с фильтрами. Пример сетки графиков на дашборде --- рис.~\ref{fig:dashboard-main}.

Scatter-графики используют равномерную подвыборку \texttt{sampleArray} (листинг~3.23) при большом числе точек на экран попадает не более \texttt{MAX\_SCATTER\_POINTS}, что сохраняет отзывчивость Recharts без потери формы облака.

\subsubsection{Карта объектов}

\texttt{ObjectsMap} загружает Yandex Maps JS API~2.1 через \texttt{yandexMapsLoader.ts}. Точки кластеризуются; лимит отображения настраивается (200--10~000, по умолчанию 2000). Точки дальше 400~км от центра города отбрасываются (листинг~3.19). Для balloon подставляется адрес из обратного геокодирования (рис.~\ref{fig:dashboard-map}).

\codelistingcaption{Отбор точек для карты (\srce{src/components/ObjectsMap.tsx})}
\begin{singleminted}{typescript}
const MAX_DISTANCE_FROM_CITY_CENTER_KM = 400;

function pointsFromRows(data, summary, coordsCache) {
    const cityCenter = resolveCityCenter(summary);
    for (const row of data) {
        const listingUrl = mapListingLinkFromRow(row);
        if (!listingUrl) continue;
        const coords = coordsFromRow(row);
        if (!coords) { missingCoords += 1; continue; }
        if (cityCenter && haversineDistanceKm(coords, cityCenter)
            > MAX_DISTANCE_FROM_CITY_CENTER_KM) {
            tooFarFiltered += 1;
            continue;
        }
        points.push({ lat: coords.lat, lng: coords.lng, listingUrl, ... });
    }
    return { points, rowsWithLink, missingCoords, tooFarFiltered };
}
\end{singleminted}
\label{lst:objects-map}

\subsubsection{Таблица, админ-панели и роли}

\texttt{FilteredListingsTable} реализует клиентскую пагинацию поверх массива \texttt{filteredData}: сортировка по числовым и текстовым колонкам, выбор видимых полей, переход на карточку объявления по колонке «Ссылка». Для строк без текстового адреса, но с координатами, хук \texttt{useReverseGeocodeAddresses} асинхронно запрашивает адрес у сервера; прогресс отображается индикатором над таблицей (рис.~\ref{fig:dashboard-table}).

\texttt{DashboardStats} агрегирует медианы цены, площади и числа комнат по текущему фильтру, дополняя их sparkline-мини-графиками. \texttt{SegmentComparison} держит два независимых набора \texttt{FilterSettings} и показывает KPI двух сегментов рынка рядом --- типичный сценарий «студии vs трёхкомнатные в пределах 5~км от центра» (рис.~\ref{fig:dashboard-segments}).

Для роли \texttt{admin} блок «Обзор рынка» заменён на \texttt{AdminServerPanel} (upload CSV, ingest) и \texttt{AdminParserPanel} (задачи парсера, прогресс, resume после капчи) --- рис.~\ref{fig:admin-server},~\ref{fig:admin-parser}.

Роль \texttt{manager} открывает \texttt{ManagerConsole}: обзор рынков и CRUD пользователей (рис.~\ref{fig:manager-markets},~\ref{fig:manager-users}). \texttt{observer} видит урезанный интерфейс без конструктора графиков и экспорта (рис.~\ref{fig:observer-dashboard}).

\codelistingcaption{Загрузка датасета по API (\srce{src/App.tsx})}
\begin{singleminted}{typescript}
const loadServerDataset = useCallback(async (cityId, opts?) => {
    if (!token) return;
    if (!opts?.silent) setLoadPhase('loading');
    try {
        const metaResp = await fetchDatasetMeta(token, cityId);
        if (metaResp) setServerMeta(metaResp.meta);
        const boot = await fetchServerBootstrap(token, cityId);
        if (!boot.ok) {
            if (caps?.canViewAdminServerPanel) {
                setLoadPhase('admin-setup');
                return;
            }
            setLoadPhase('server-empty');
            return;
        }
        applyDataset(boot.data.rows, boot.data.summary, null);
    } catch {
        setToast('Сервер недоступен — запустите npm run dev:api');
        setLoadPhase('admin-setup');
    }
}, [token, applyDataset, caps]);
\end{singleminted}
\label{lst:load-server}

\codelistingcaption{Фрагмент панели фильтров (\srce{src/components/FilterPanel.tsx})}
\begin{singleminted}{typescript}
const filtersDirty = useMemo(
    () => isAnyFilterDirty(filters, baselineFilters),
    [filters, baselineFilters]
);

const resetFilters = () => {
    onFilterChange({
        ...baselineFilters,
        price: { ...baselineFilters.price },
        area: { ...baselineFilters.area },
        rooms: [...baselineFilters.rooms],
        excludeFirstFloor: baselineFilters.excludeFirstFloor,
        excludeLastFloor: baselineFilters.excludeLastFloor,
        radiusKm: baselineFilters.radiusKm,
    });
};

// onChange слайдера расстояния синхронизирует radiusKm с max диапазона
onFilterChange({ ...filters, distanceKm: next, radiusKm: next.max });
\end{singleminted}
\label{lst:filter-panel}

% ---------------------------------------------------------------------------
\subsection{Алгоритмы визуализации и производительности}
% ---------------------------------------------------------------------------

\subsubsection{Биннинг гистограмм}

Функция \texttt{buildHistogram} в \texttt{DynamicChartGrid.tsx} (листинг~3.20) делит значения на 20 интервалов (\texttt{HIST\_BINS = 20}). По умолчанию границы bins берутся по 2-му и 98-му перцентилям (\texttt{quantile}), чтобы выбросы не сжимали основную массу столбцов; крайние bin'ы аккумулируют «хвосты». В подписи указываются медиана и IQR. Пользователь может увеличить интервал (zoom), тогда пересчёт выполняется только по значениям внутри выбранного диапазона.

\codelistingcaption{Гистограмма с перцентильными границами (\srce{src/components/DynamicChartGrid.tsx})}
\begin{singleminted}{typescript}
const HIST_BINS = 20;

function buildHistogram(values: number[], zoom?) {
    const min = Math.min(...values);
    const max = Math.max(...values);
    const qLo = quantile(values, 0.02);
    const qHi = quantile(values, 0.98);
    const useCore = (qHi - qLo) > 0 && (qHi - qLo) < (max - min) * 0.95;
    const binStart = useCore ? qLo : min;
    const binEnd = useCore ? qHi : max;
    const binSize = (binEnd - binStart || 1) / HIST_BINS;
    const distribution = Array.from({ length: HIST_BINS }, (_, index) => ({
        name: formatNumber(binStart + index * binSize),
        count: 0,
        range: getBinRange(binStart, binSize, index, HIST_BINS, binEnd),
    }));
    for (const value of values) {
        let binIndex = value < binStart ? 0
            : value > binEnd ? HIST_BINS - 1
            : Math.floor((value - binStart) / binSize);
        distribution[binIndex].count += 1;
    }
    return { distribution, analysis: `Медиана: ${median(values)}`, ... };
}
\end{singleminted}
\label{lst:histogram}

\subsubsection{Подвыборка scatter и кэширование}

Для scatter-графиков функция \texttt{sampleArray} ограничивает число точек константой \texttt{MAX\_SCATTER\_POINTS} (1200), сохраняя равномерное покрытие диапазона --- иначе Recharts тормозит на десятках тысяч объектов. Аналогично \texttt{samplePointsWithinLimit} на карте равномерно прореживает маркеры при превышении пользовательского лимита.

\subsubsection{Обратное геокодирование на клиенте}

Клиент не вызывает Yandex для каждой строки: \texttt{useReverseGeocodeAddresses} отправляет ключи на \srce{GET /api/geocode/reverse} и \srce{POST /api/geocode/progress}, используя серверный SQLite-кэш.

\subsubsection{Инструментирование и вспомогательные модули}

В режиме \texttt{import.meta.env.DEV} компоненты логируют время фильтрации и построения гистограмм при превышении порога 4~мс. Ниже приведены вспомогательные фрагменты: npm-скрипты, проверка схемы QA, прореживание scatter, one-hot в ETL и загрузчик карт.

\codelistingcaption{Скрипты npm (\srce{package.json})}
\begin{singleminted}{json}
{
  "scripts": {
    "dev": "vite",
    "dev:api": "tsx server/index.ts",
    "dev:all": "npm run dev:api & npm run dev",
    "build": "vite build",
    "check:parser-configs": "tsx server/validateParserConfigs.ts"
  }
}
\end{singleminted}
\label{lst:package-json}

Команда \texttt{dev:all} поднимает API и клиент одновременно; \texttt{check:parser-configs} вызывается перед длительным сбором данных.

\codelistingcaption{Валидация схемы в QA (\srce{parser/pipeline/qa_report.py})}
\begin{singleminted}{python}
def validate_schema(df, schema):
    errors, warnings = [], []
    required_core = schema.get("required_core", [])
    for col in required_core:
        if col not in df.columns:
            errors.append(f"Обязательная колонка отсутствует: {col}")
        elif df[col].isna().mean() > 0.15:
            warnings.append(f"Доля пропусков в `{col}` > 15%.")
    return errors, warnings
\end{singleminted}
\label{lst:qa-schema}

Функция \texttt{validate\_schema} отделяет ошибки (нет обязательной колонки) от предупреждений (много пропусков).

\codelistingcaption{Подвыборка для scatter (\srce{src/utils/sample.ts})}
\begin{singleminted}{typescript}
export const MAX_SCATTER_POINTS = 1200;

export function sampleArray<T>(items: T[], limit: number): T[] {
    if (items.length <= limit) return items;
    const lastIndex = items.length - 1;
    return Array.from({ length: limit }, (_, idx) => {
        const sourceIndex = Math.floor((idx * lastIndex) / (limit - 1));
        return items[sourceIndex];
    });
}
\end{singleminted}
\label{lst:sample-array}

Равномерная подвыборка сохраняет форму облака точек при большом числе объявлений.

\codelistingcaption{One-hot кодирование в ETL (\srce{parser/pipeline/etl.py})}
\begin{singleminted}{python}
def apply_one_hot(out, row, output_columns):
    for prefix, (src_col, classify) in CATEGORY_CLASSIFIERS.items():
        category = classify(row.get(src_col))
        cols = one_hot_columns(prefix, output_columns)
        for col in cols:
            out[col] = False
        target = f"{prefix}_{category}"
        if target in cols:
            out[target] = True
\end{singleminted}
\label{lst:one-hot}

Категориальные признаки карточки превращаются в бинарные столбцы для фильтров дашборда.

\codelistingcaption{Загрузка Yandex Maps API (\srce{src/utils/yandexMapsLoader.ts})}
\begin{singleminted}{typescript}
export function getYandexMapsApiKey(): string {
    return (import.meta.env.VITE_YANDEX_MAPS_API_KEY ?? '').trim();
}

export function loadYandexMaps(apiKey: string): Promise<Ymaps21Global> {
    const key = apiKey.trim();
    if (!key) throw new Error('VITE_YANDEX_MAPS_API_KEY не задан');
    return withTimeout(
        loadYmaps21Script(key, DEFAULT_LOAD_TIMEOUT_MS),
        DEFAULT_LOAD_TIMEOUT_MS,
        'Yandex Maps'
    );
}
\end{singleminted}
\label{lst:yandex-maps-loader}

Загрузчик проверяет наличие ключа из \texttt{.env} и ограничивает время ожидания скрипта API.

\subsubsection{Согласованность контрактов между слоями}

Сквозная согласованность данных обеспечивается несколькими «контрактами». JSON \texttt{schema\_contract.json} задаёт список колонок processed CSV и используется в ETL, QA и при построении \texttt{DataSummary} на сервере. Модуль \texttt{shared/dashboard.ts} описывает типы \texttt{DataRow}, \texttt{FilterSettings}, \texttt{DataSummary}; клиент импортирует их через алиас \texttt{@/types}, сервер --- напрямую из \texttt{../shared/}. Профили городов в \texttt{cities.ts} дублируют координаты центра из Python-конфига, чтобы расстояние до центра, рассчитанное в ETL, совпадало с радиусом фильтра на карте.

Таблица~\ref{tab:layer-contracts} суммирует ключевые точки стыка слоёв.

\begin{table}[htbp]
\centering
\footnotesize
\setlength{\tabcolsep}{3pt}
\renewcommand{\arraystretch}{1.15}
\begin{tabular}{|L{2.9cm}|L{5.5cm}|L{5.0cm}|}
\hline
\textbf{Стык} & \textbf{Артефакт} & \textbf{Назначение} \\
\hline
Parser $\rightarrow$ Server & \texttt{processed\_apartment\_\newline data.csv} & Ingest через admin API \\
\hline
Server $\rightarrow$ Client & \texttt{GET /api/dataset/\newline bootstrap} & JSON датасета и summary \\
\hline
Shared фильтры & \texttt{rowPassesFilters\newline Shared} & Единый срез таблицы и API \\
\hline
Runtime парсера & \texttt{.runtime\_status.\newline json} & Прогресс в AdminParserPanel \\
\hline
Геокодер & \texttt{reverse\_geocode\_cache} & Кэш адресов в SQLite \\
\hline
\end{tabular}
\caption{Ключевые контракты между слоями программного комплекса}
\label{tab:layer-contracts}
\end{table}

Такой подход позволяет заменить или масштабировать один слой --- например, перенести хранение датасета в PostgreSQL --- без переписывания фильтров и клиентских компонентов, если сохранить форму JSON bootstrap и summary.

\subsubsection{Модульные тесты и smoke-проверки парсера}

В каталоге \texttt{parser/pipeline/} собраны точечные тесты: \texttt{test\_captcha\_detection.py} проверяет чистую логику \texttt{should\_treat\_as\_captcha\_page} без запуска браузера; \texttt{test\_config\_presets.py} --- корректность пресетов задержек; \texttt{test\_write\_rows.py} --- атомарность \texttt{DetailsStore.write\_rows}. Скрипт \texttt{smoke\_test.py} выполняет короткий прогон links/details на малом лимите ссылок и используется перед демонстрацией системы.

На стороне сервера команда \texttt{npm run check:parser-configs} обходит все города из \texttt{CITY\_PROFILES} и выводит ошибки конфигурации в консоль CI или локальной разработки. Таким образом, регрессии в путях к файлам или опечатки в URL сегментов обнаруживаются до запуска многочасового сбора.

% ---------------------------------------------------------------------------
\subsection{Развёртывание и эксплуатация}
% ---------------------------------------------------------------------------

Помимо локальной разработки, комплекс подготовлен к выкладке в сеть: статический фронтенд, API за reverse-proxy и автозапуск через systemd. Конфигурация и скрипты собраны в каталоге \texttt{deploy/}.

\subsubsection{Режимы запуска}

В ходе работы использовались три режима (табл.~\ref{tab:deploy-modes}).

\begin{table}[htbp]
\centering
\small
\setlength{\tabcolsep}{3pt}
\begin{tabular}{|L{2.6cm}|L{3.4cm}|L{7.4cm}|}
\hline
\textbf{Режим} & \textbf{Команда} & \textbf{Назначение} \\
\hline
Разработка & \texttt{npm run dev:all} & Vite на порту 5173, API на 3001; прокси \texttt{/api} в \texttt{vite.config.ts} \\
\hline
Демонстрация по URL & \texttt{npm run dev:tunnel} & Production-сборка + \texttt{vite preview} на 3000, API на 3001, HTTPS-туннель (localtunnel по умолчанию) \\
\hline
Production на VPS & \texttt{deploy/setup.sh} & nginx отдаёт \texttt{dist/}, проксирует \texttt{/api} на Node.js, сервис systemd \\
\hline
\end{tabular}
\caption{Режимы запуска программного комплекса}
\label{tab:deploy-modes}
\end{table}

Команда \texttt{dev:all} применялась при отладке UI и API на одной машине. Для показа дашборда руководителю или проверки с мобильного устройства использовался сценарий \texttt{dev:tunnel}: скрипт \texttt{scripts/dev-with-tunnel.sh} собирает фронт (\texttt{npm run build}), поднимает \texttt{vite preview} и Express, затем пробрасывает порт наружу через localtunnel, Cloudflare Quick Tunnel или ngrok (переменная \texttt{TUNNEL\_PROVIDER}). Режим \texttt{TUNNEL\_SERVE=preview} выбран осознанно: dev-сервер Vite через туннель часто отдаёт 502 и белый экран, тогда как preview production-бандла работает стабильнее.

\subsubsection{Production-сборка фронтенда}

Сборка выполняется одной командой \texttt{npm run build}. Vite компилирует React-приложение в каталог \texttt{dist/}: минифицированный JS/CSS, \texttt{index.html}, ассеты. Переменные \texttt{VITE\_*} (\texttt{VITE\_YANDEX\_MAPS\_API\_KEY} и при необходимости \texttt{VITE\_API\_BASE\_URL}) \emph{вшиваются на этапе сборки}; после смены ключа карт требуется повторный \texttt{build}. Серверная часть (\texttt{tsx server/index.ts}) не компилируется отдельно --- на VPS запускается через \texttt{npm run start:api} с \texttt{NODE\_ENV=production}, что отключает отладочные эндпоинты.

\subsubsection{Схема развёртывания на VPS}

Итоговая схема на выделенном сервере (VPS Beget, Ubuntu~24.04) представлена на рис.~\ref{fig:deploy-vps}. Пользователь обращается к nginx по HTTP (порт~80); nginx отдаёт файлы из \texttt{dist/} и для путей \texttt{/api/} проксирует запросы на \texttt{127.0.0.1:3001}, где слушает Express. SQLite и JSON-датасеты остаются в \texttt{server/data/} на том же хосте; Python-парсер по-прежнему запускается subprocess из API при действиях администратора.

\begin{figure}[htbp]
\centering
\begin{tikzpicture}[
    comp/.style={draw, rounded corners=2pt, align=center, minimum height=11mm, inner sep=4pt, font=\small},
    arr/.style={-{Stealth[length=2.2mm]}, semithick},
    lbl/.style={font=\scriptsize, fill=white, inner sep=1.5pt}
]
\node[comp, fill=blue!8] (user) {Браузер\\пользователя};
\node[comp, fill=orange!10, below=10mm of user] (nginx) {nginx :80\\статика + proxy};
\node[comp, fill=gray!10, below left=9mm and -8mm of nginx] (dist) {\texttt{dist/}\\SPA};
\node[comp, fill=gray!10, below right=9mm and -8mm of nginx] (api) {Node.js API\\127.0.0.1:3001};
\node[comp, fill=gray!10, below=8mm of api] (data) {SQLite + JSON\\ \texttt{server/data/}};
\draw[arr] (user) -- node[right, lbl] {HTTP} (nginx);
\draw[arr] (nginx) -| node[near start, above, lbl] {\texttt{/}} (dist);
\draw[arr] (nginx) -| node[near start, above, lbl] {\texttt{/api}} (api);
\draw[arr] (api) -- (data);
\end{tikzpicture}
\caption{Развёртывание на VPS: nginx, статика и API на одном сервере}
\label{fig:deploy-vps}
\end{figure}

Конфигурация nginx (\texttt{deploy/nginx-dashboard.conf}) задаёт \texttt{root} на \texttt{dist/}, \texttt{try\_files} для SPA-маршрутизации и \texttt{proxy\_pass} на backend; лимит тела запроса \texttt{client\_max\_body\_size 25m} согласован с загрузкой CSV через admin API.

\subsubsection{Автоматизация установки}

Первичная настройка сервера выполняется скриптом \texttt{deploy/setup.sh} (запуск с \texttt{sudo} из корня репозитория). Скрипт последовательно:
\begin{enumerate}
    \item проверяет наличие \texttt{.env} (шаблон --- \texttt{deploy/env.production.example});
    \item устанавливает Node.js~20 (через NodeSource), nginx и \texttt{build-essential} (нативная сборка \texttt{better-sqlite3});
    \item выполняет \texttt{npm install} и \texttt{npm run build} с подгруженными переменными окружения;
    \item размещает конфиг nginx в \texttt{/etc/nginx/sites-available/}, включает сайт, перезагружает nginx;
    \item регистрирует unit \texttt{real-estate-dashboard.service} в systemd (\texttt{Restart=on-failure}, \texttt{EnvironmentFile=.env}) и запускает API.
\end{enumerate}

Перед установкой на сервер клонируется репозиторий (например, в \texttt{/var/www/real-estate-analytics-dashboard}), создаётся \texttt{.env} с \texttt{JWT\_SECRET}, \texttt{YANDEX\_GEOCODER\_API\_KEY} и \texttt{VITE\_YANDEX\_MAPS\_API\_KEY}. Переменная \texttt{SERVER\_NAME} передаётся в \texttt{setup.sh} и подставляется в \texttt{server\_name} nginx. \texttt{VITE\_API\_BASE\_URL} не задаётся: клиент и API доступны с одного IP/домена, запросы идут на относительный путь \texttt{/api}.

\subsubsection{Ключи Яндекса и проверка после деплоя}

После выкладки в кабинете разработчика Яндекса (рис.~\ref{fig:yandex-api}) для JavaScript API карт в список Referer добавляется адрес сервера, например \texttt{http://<IP>/*}. Без этого карта на production остаётся серой, хотя локально всё работает.

Проверка на сервере:
\begin{itemize}
    \item \texttt{curl http://127.0.0.1/api/health} --- ответ JSON от API;
    \item \texttt{curl -I http://127.0.0.1/} --- HTTP~200 и отдача \texttt{index.html};
    \item \texttt{systemctl status real-estate-dashboard} --- активный unit systemd.
\end{itemize}

В браузере открывается главная страница дашборда; вход выполняется под seed-учётными записями (\texttt{observer}, \texttt{analyst}, \texttt{admin}). После первого входа рекомендуется сменить пароли.

\subsubsection{Обновление версии}

Обновление на VPS сводится к \texttt{git pull}, повторному \texttt{npm install} при изменении зависимостей, \texttt{npm run build} (с актуальным \texttt{.env}) и \texttt{systemctl restart real-estate-dashboard}. Статика в \texttt{dist/} подхватывается nginx без перезапуска; перезапуск нужен только API. Логи API: \texttt{journalctl -u real-estate-dashboard -f}.

% ---------------------------------------------------------------------------
\subsection{Выводы по третьему разделу}
% ---------------------------------------------------------------------------

В третьем разделе описана программная реализация комплекса на уровне исходного кода:
\begin{enumerate}
    \item Python-пайплайн на Selenium с поэтапным сбором ссылок и карточек, ETL по контракту схемы, QA-отчётами и возобновлением после SmartCaptcha через JSON-файлы \texttt{.runtime\_status.json} и \texttt{.runtime\_control.json};
    \item Express-сервер с JWT-авторизацией по ролям, ingest JSON-датасетов по городам, spawn Python-парсера с bootstrap venv и многоуровневым обратным геокодированием с кэшем SQLite;
    \item общий слой \texttt{shared/}, исключающий расхождение правил фильтрации, профилей городов и слияния строк между клиентом и API;
    \item React-клиент с отложенной фильтрацией (\texttt{useDeferredValue}), чанковым CSV-импортом, настраиваемыми графиками Recharts с перцентильным биннингом и картой на API Яндекса с кластеризацией и лимитом точек;
    \item сценарии развёртывания: локальная разработка, демонстрация через HTTPS-туннель и production на VPS (nginx + systemd) с автоматизированной установкой из каталога \texttt{deploy/}.
\end{enumerate}

Каждый слой имеет чёткую точку входа, конфигурацию и наблюдаемый статус выполнения.


\section{Тестирование и результаты}

В четвёртом разделе описывается проверка разработанного программного комплекса на реальных данных Яндекс.Недвижимость, оценка качества сбора и подготовки датасетов, а также пример прикладного анализа рынка через дашборд.

\subsection{Методика проверки}

Проверка проводилась без отдельного фреймворка автоматизированных UI-тестов: для учебного проекта достаточно сочетания контроля качества данных (QA-пайплайн), ручных сценариев по ролям и наблюдения за поведением системы на реальных выборках Москвы и Саратова.

\textbf{Уровни проверки:}
\begin{enumerate}
    \item \textbf{Данные парсера} --- стадии links, details, ETL; отчёт QA по контракту схемы; сверка числа строк после ingest на сервер.
    \item \textbf{Серверное API} --- авторизация, bootstrap датасета, фильтрация и экспорт CSV; запуск и мониторинг задач парсера из админ-панели.
    \item \textbf{Клиентский дашборд} --- загрузка серверного снимка и локального CSV; фильтры, KPI, графики, карта, таблица; разграничение возможностей по ролям.
    \item \textbf{Развёртывание} --- production-сборка, доступ по публичному URL (туннель или VPS), health-check API, работа карт с Referer production-хоста.
    \item \textbf{Интеграции} --- JavaScript API Яндекс.Карт на клиенте, API Геокодера и кэш SQLite на сервере (см.~рис.~\ref{fig:yandex-api}).
\end{enumerate}

\textbf{Учётные записи для сценариев} (seed при первом запуске сервера): \texttt{observer}, \texttt{analyst}, \texttt{admin}, \texttt{manager} с паролями, заданными в документации проекта. Каждая роль проходила типовой сценарий из табл.~\ref{tab:roles} (раздел~2).

\textbf{Критерии успешности:}
\begin{itemize}
    \item датасет города загружается в браузер без ошибок и отображает тысячи строк (Москва) или демонстрационную выборку (Саратов);
    \item фильтры меняют KPI и графики без заметной блокировки интерфейса;
    \item экспорт отфильтрованного CSV совпадает с таблицей на экране;
    \item администратор может запустить стадию парсера и увидеть прогресс; после ETL --- выполнить ingest;
    \item QA не фиксирует отсутствие обязательных колонок в processed CSV.
\end{itemize}

\subsection{Проверка пайплайна сбора и качества данных}

\subsubsection{Сценарий сбора для Саратова}

Для Саратова выполнен полный цикл: \texttt{links} $\rightarrow$ \texttt{details} $\rightarrow$ \texttt{etl} с ограничением \texttt{target\_listings=100}. По логам runtime-статуса: собрано порядка 1174 ссылок, обработано 95 карточек, после ETL --- 95 строк в \texttt{processed\_apartment\_data.csv}. Стадия завершилась без ошибок парсера; ingest добавил записи в серверный JSON-датасет с дедупликацией по ID объявления.

При повторном append часть строк была пропущена как дубликаты (в meta зафиксированы \texttt{skippedDuplicates}), что подтверждает работу слияния при загрузке.

\subsubsection{Сценарий для Москвы}

Московский датасет формировался несколькими партиями append из processed CSV. Итоговый серверный снимок содержит \textbf{12\,549} объявлений (обновление от 25.05.2026). В meta сохранены две загрузки: первая партия --- 5341 строка, вторая --- 7208 добавленных при 3217 пропущенных дубликатах и 16 отфильтрованных как слишком удалённых от центра города.

\subsubsection{Отчёт QA}

Модуль \texttt{pipeline/qa\_report.py} проверяет наличие \texttt{required\_core} колонок (цена, комнаты, площадь, координаты), долю пропусков (предупреждение при $>$15\%), дубликаты ссылок. Для сданных датасетов критических ошибок схемы не выявлено; предупреждения возможны по отдельным полям с высокой долей пропусков (например, текстовый адрес, если координаты есть, но reverse geocode ещё не применён).

\subsection{Результаты на данных Москвы и Саратова}

Сводные показатели серверных снимков приведены в табл.~\ref{tab:results-cities}.

\begin{table}[htbp]
\centering
\small
\begin{tabular}{|L{2.4cm}|L{2.6cm}|L{2.8cm}|L{2.8cm}|L{2.8cm}|}
\hline
\textbf{Город} & \textbf{Число объявлений} & \textbf{Медиана цены, млн руб.} & \textbf{Медиана площади, м$^2$} & \textbf{Медиана расст. до центра, км} \\
\hline
Москва & 12\,549 & 24,0 & 55,2 & 11,0 \\
\hline
Саратов & 75 & 4,1 & 39,8 & 7,4 \\
\hline
\end{tabular}
\caption{Сводные характеристики серверных датасетов (по состоянию на 25.05.2026)}
\label{tab:results-cities}
\end{table}

\textbf{Москва} --- основная выборка для нагрузочной проверки клиента: bootstrap JSON $\approx$47~МБ, фильтрация и перерисовка графиков на десятках тысяч строк остаются приемлемыми за счёт \texttt{useMemo} и \texttt{useDeferredValue} (рис.~\ref{fig:dashboard-main}). Карта с лимитом 2000 точек по умолчанию показывает кластеры по агломерации (рис.~\ref{fig:dashboard-map}); увеличение лимита до нескольких тысяч возможно, но увеличивает время отрисовки.

\textbf{Саратов} --- демонстрационный регион с меньшим объёмом данных; удобен для быстрой проверки полного цикла парсера и ingest без длительного ожидания details.

\subsubsection{Проверка развёртывания}

Помимо локального \texttt{npm run dev:all}, система проверялась в двух сетевых конфигурациях. \textbf{Демонстрационный туннель} (\texttt{npm run dev:tunnel}): production-бандл и API пробрасываются наружу через localtunnel; подтверждены вход по JWT, загрузка bootstrap Москвы и отображение карты при корректном Referer. \textbf{VPS} (Ubuntu на Beget): выполнена установка скриптом \texttt{deploy/setup.sh} --- nginx на порту~80, unit \texttt{real-estate-dashboard}, статика из \texttt{dist/}, проксирование \texttt{/api} на \texttt{127.0.0.1:3001}. После деплоя успешно проходят \texttt{GET /api/health}, открытие главной страницы и сценарии analyst/admin из табл.~\ref{tab:roles-test}. Схема production-размещения --- рис.~\ref{fig:deploy-vps} (раздел~3).

\subsection{Проверка сценариев по ролям}

Результаты ручной проверки согласованы с матрицей доступа (табл.~\ref{tab:roles-test}).

\begin{table}[htbp]
\centering
\small
\begin{tabular}{|L{2.2cm}|L{4.2cm}|L{7.0cm}|}
\hline
\textbf{Роль} & \textbf{Сценарий} & \textbf{Результат} \\
\hline
observer & Вход, Москва, фильтры, KPI, графики, карта & Доступен только серверный датасет; таблица и экспорт недоступны; упрощённый список объявлений отображается \\
\hline
analyst & Сервер + свой CSV, конструктор графиков, экспорт & Переключение источника работает; CSV до 50\,000 строк загружается чанками; экспорт CSV/XLSX совпадает с фильтром \\
\hline
admin & Upload, ingest, запуск парсера, прогресс & Ingest processed CSV обновляет meta; parser job показывает стадию и процент; resume после капчи принимается через API \\
\hline
manager & Вкладки «Рынки» и «Пользователи» & Сводка по городам совпадает с meta; создание пользователя и смена роли сохраняются в SQLite \\
\hline
\end{tabular}
\caption{Результаты проверки типовых сценариев по ролям}
\label{tab:roles-test}
\end{table}

Проверка роли \texttt{observer} подтвердила доступ только к серверному датасету и отсутствие таблицы с экспортом (рис.~\ref{fig:observer-dashboard}). Для \texttt{analyst} работают переключение источника, конструктор графиков и выгрузка --- на рис.~\ref{fig:dashboard-table}. Администратор успешно выполнял ingest (рис.~\ref{fig:admin-server}) и отслеживал задачи парсера (рис.~\ref{fig:admin-parser}). Руководитель видит сводку по рынкам (рис.~\ref{fig:manager-markets}) и управляет учётными записями (рис.~\ref{fig:manager-users}).

\textbf{Авторизация:} неверный пароль возвращает HTTP 401; без токена защищённые маршруты недоступны. JWT действует 12~часов, после истечения требуется повторный вход.

\textbf{Согласованность фильтров:} при одинаковых параметрах query клиентская таблица, KPI и \texttt{GET /api/listings} дают одинаковое число строк (проверено на выборке «двушки, цена до 15~млн, расстояние до 10~км» для Москвы).

\subsection{Пример прикладного анализа}

Рассмотрим сегмент вторичного рынка двухкомнатных квартир в Москве в радиусе 10~км от центра с ценой до 15~млн руб. (типичный сценарий из раздела~2). Состояние дашборда после применения фильтров показано на рис.~\ref{fig:analysis-filters}.

\screenshot{analysis-filters}{analysis-filters}{Прикладной срез: двухкомнатные квартиры, цена до 15~млн руб., расстояние до центра до 10~км (Москва)}

После применения фильтров в выборке остаётся порядка нескольких сотен объявлений (точное число зависит от актуальности снимка). Медиана цены оказывается ниже общегородской медианы (24~млн), медиана площади --- около 45--55~м$^2$. На гистограмме цены (рис.~\ref{fig:analysis-charts}) виден левый сдвиг распределения; scatter «цена vs площадь» показывает положительную связь с разбросом на уровне старых панельных серий.

\screenshot{analysis-charts}{analysis-charts}{Гистограмма цены и scatter «цена~--- площадь» для отфильтрованного сегмента}

Блок \textbf{«Сравнение сегментов»} позволяет рядом поставить, например, «новостройки» (фильтр по году постройки $\geq$ 2015) и оставшуюся вторичку (рис.~\ref{fig:analysis-segments}): медиана цены за м² у новостроек, как правило, выше, что согласуется с общими представлениями о премиальности первичного рынка.

\screenshot{analysis-segments}{analysis-segments}{Сравнение новостроек (год $\geq$ 2015) и вторичного рынка в том же географическом срезе}

На \textbf{карте} объекты сегмента группируются в кольце вокруг центра (рис.~\ref{fig:dashboard-map}); балуны содержат цену и ссылку на карточку. Для строк без адреса выполнялось обратное геокодирование: часть записей получила текстовый адрес через кэш SQLite после нескольких запросов к API.

\textbf{Экспорт:} отфильтрованная таблица выгружена в CSV и сверена с экраном --- расхождений по числу строк нет.

\subsection{Производительность и ограничения}

При работе с московским снимком (12\,549 строк) первичная загрузка bootstrap занимает несколько секунд на локальном сервере; последующие изменения фильтров не вызывают повторных запросов к API. Импорт CSV аналогичного размера на клиенте выполняется с паузами между чанками --- интерфейс остаётся отзывчивым.

Выявленные \textbf{ограничения}:
\begin{itemize}
    \item парсер чувствителен к капче Яндекса --- требуется ручное вмешательство или headless$\rightarrow$окно;
    \item не все объявления содержат полный адрес; геокодирование ограничено квотами API;
    \item SQLite и JSON-файлы рассчитаны на локальное/учебное развёртывание, не на высокую конкурентную нагрузку;
    \item саратовская выборка мала для устойчивых выводов о структуре рынка --- годится для отладки, не для статистической репрезентативности.
\end{itemize}

\subsection{Выводы по четвёртому разделу}

Проведённая проверка показала, что программный комплекс выполняет заявленные функции:
\begin{enumerate}
    \item пайплайн Selenium собирает и нормализует данные; QA и ingest обеспечивают контролируемое попадание строк на сервер;
    \item серверные датасеты Москвы (12\,549 объявлений) и Саратова (75 объявлений) доступны через API и дашборд;
    \item ролевая модель ограничивает операции записи и экспорт в соответствии с проектом;
    \item интерактивный анализ (фильтры, графики, карта, сравнение сегментов) применим к реальным данным и поддерживает прикладные срезы рынка;
    \item развёртывание на VPS и демонстрация через HTTPS-туннель подтверждают пригодность комплекса для показа вне локальной машины разработчика.
\end{enumerate}

Система пригодна для учебной и исследовательской работы с объявлениями Яндекс.Недвижимость; для промышленной эксплуатации потребовались бы масштабирование хранилища, автоматизация UI-тестов и более устойчивый сбор данных в условиях anti-bot защиты.


\conclusion

В работе разработан программный комплекс для анализа рынка недвижимости на данных платформы Яндекс.Недвижимость: Python-пайплайн сбора, серверное хранение датасетов по городам и интерактивный веб-дашборд с разграничением прав доступа.

В первом разделе рассмотрены теоретические основы сбора данных из динамических веб-источников, предобработки табличных массивов и построения аналитических панелей. Обоснован переход к автоматизации браузера (Selenium) для сайтов с JavaScript-контентом.

Во втором разделе спроектирована архитектура системы: клиент (React SPA), сервер API (Express, SQLite, JSON-датасеты), модуль парсера и общий слой правил \texttt{shared/}. Описаны модель данных, REST-интерфейс, роли пользователей и подключение API Яндекс.Карт и Геокодера.

В третьем разделе приведена программная реализация: поэтапный парсер с ETL и обработкой капчи, ingest на сервер, единая фильтрация, клиентские визуализации на Recharts и карта объектов.

В четвёртом разделе выполнена проверка на данных Москвы (12\,549 объявлений) и Саратова (75 объявлений). Подтверждена работоспособность сценариев для ролей observer, analyst, admin и manager; показан пример прикладного среза рынка через фильтры дашборда. Система развёрнута на VPS (nginx, systemd) и проверена через публичный HTTPS-туннель.

Практический результат --- работающее приложение полного цикла «сбор $\rightarrow$ хранение $\rightarrow$ интерактивный анализ», пригодное для учебных и исследовательских задач по локальным рынкам жилой недвижимости.

\bibliographystyle{ugost2003}
\bibliography{thesis}

\appendix

% Полные тексты исходных файлов (дублируют фрагменты листингов раздела~3).
% VerbatimInput вместо minted: иначе pdfTeX исчерпывает пул строк (~10k строк кода).

\newcommand{\appendixfile}[2]{%
  \section{#1}%
  \vspace{0.4\baselineskip}%
  \noindent{\small\srce{#2}}\par
  \vspace{0.6\baselineskip}%
  \begin{spacing}{0.92}%
  \begingroup
  \fvset{breaklines=true,breakanywhere=true,fontsize=\scriptsize,formatcom=\scriptsize\ttfamily}%
  \VerbatimInput{#2}%
  \endgroup
  \par\vspace{1.2\baselineskip}%
  \end{spacing}%
}

\appendixfile{Конфигурация пайплайна (config.py)}{../parser/pipeline/config.py}
\appendixfile{Стадия links (run\_links.py)}{../parser/pipeline/run_links.py}
\appendixfile{Хранилище details CSV (details\_store.py)}{../parser/pipeline/details_store.py}
\appendixfile{Парсер карточек объявлений (yandex\_realty\_details\_parser.py)}{../parser/yandex_realty_details_parser.py}
\appendixfile{Модуль ETL (etl.py)}{../parser/pipeline/etl.py}
\appendixfile{Оркестратор pipeline (run\_pipeline.py)}{../parser/pipeline/run_pipeline.py}
\appendixfile{Runtime-статус и resume (runtime\_status.py)}{../parser/pipeline/runtime_status.py}
\appendixfile{Отчёт QA (qa\_report.py)}{../parser/pipeline/qa_report.py}

\appendixfile{Точка входа Express API (index.ts)}{../server/index.ts}
\appendixfile{Хранение и ingest датасетов (datasetStore.ts)}{../server/datasetStore.ts}
\appendixfile{Запуск Python-парсера (parserRunner.ts)}{../server/parserRunner.ts}
\appendixfile{Обратное геокодирование (reverseGeocoder.ts)}{../server/reverseGeocoder.ts}

\appendixfile{Профили городов (cities.ts)}{../shared/cities.ts}
\appendixfile{Фильтрация строк (filters.ts)}{../shared/filters.ts}
\appendixfile{Слияние строк датасета (mergeDatasetRows.ts)}{../shared/mergeDatasetRows.ts}

\appendixfile{Корневой компонент приложения (App.tsx)}{../src/App.tsx}
\appendixfile{Импорт CSV на клиенте (csvImport.ts)}{../src/utils/csvImport.ts}
\appendixfile{Панель фильтров (FilterPanel.tsx)}{../src/components/FilterPanel.tsx}
\appendixfile{Сетка графиков (DynamicChartGrid.tsx)}{../src/components/DynamicChartGrid.tsx}
\appendixfile{Карта объектов (ObjectsMap.tsx)}{../src/components/ObjectsMap.tsx}
\appendixfile{Подвыборка для scatter (sample.ts)}{../src/utils/sample.ts}
\appendixfile{Загрузчик Yandex Maps API (yandexMapsLoader.ts)}{../src/utils/yandexMapsLoader.ts}

\appendixfile{Манифест npm (package.json)}{../package.json}


\end{document}
